\section{Optimizers near the singular endpoint}
\label{sec:singular-endpoint}

In this section, we prove \Cref{thm:endpoint-optimizers} in a more precise form.
Along the Lubetzky--Zhao boundary, the supporting line to the graph of
$x\mapsto J_{\pc(r)}(x^{1/d})$ touches it at $x=r^d$ and
$x=\sm(r)^d$.  At the exceptional density $r_*=(d-1)/d$, these two
points coincide at $x=r_*^d$, so the nonexceptional reduction no longer
applies.  Following the strategy in
\Cref{sec:proof-strategy}, we construct an analytic rank-one bipodal family
and establish its global optimality and uniqueness along the associated
parameter curve approaching $(\pc(r_*),r_*)$.

Fix a finite simple $d$-regular graph $H$ with $d\ge2$,
and write $m:=|E(H)|$ and $v:=|V(H)|=2m/d$.  Set
\begin{equation*}
  r_*=\frac{d-1}{d},
  \qquad
  p_*=\pc(r_*)
  =\frac{d-1}{(d-1)+\exp(d/(d-1))},
  \qquad
  u_*:=\sqrt{r_*},
  \qquad
  P_*:=(p_*,r_*).
\end{equation*}
We write $\lambda$ for Lebesgue measure on $[0,1]$, and also for its
product measure on $[0,1]^2$.
In this section, a kernel means a measurable real-valued function of two
variables; unlike a graphon, it need not be symmetric or take values in
$[0,1]$.
For a measurable function $g:[0,1]\to\mathbb R$, write $g\otimes g$ for
the rank-one kernel $(g\otimes g)(x,y):=g(x)g(y)$.
The following theorem is the main result of this section.
The remainder of this section is devoted to proving this theorem.

\begin{theorem}[Optimizers near the singular endpoint]
\label{thm:singular-endpoint}
There exist $h_0>0$ and real-analytic functions
$p_h,r_h,u_h,\alpha_h$, defined for $|h|<h_0$, such that
\[
  (p_0,r_0)=P_*,
  \qquad
  u_0=u_*,
  \qquad
  \alpha_0=\frac12.
\]
For every $h\in(0,h_0)$, we have $p_h<\pc(r_h)$,
and for any measurable set $A_h\subset[0,1]$
with $\lambda(A_h)=\alpha_h$, set
\[
  s_h:=u_h-h,
  \qquad
  t_h:=u_h+h,
  \qquad
  f_h:=s_h\1_{A_h}+t_h\1_{A_h^c},
  \qquad
  W_h:=f_h\otimes f_h.
\]
Then $W_h$ is a nonconstant rank-one bipodal graphon satisfying $t(H,W_h)=r_h^m$
and is the unique minimizer of \eqref{eq:graphon-variational-problem}
at $(p_h,r_h)$, up to relabeling.  Moreover,
\begin{equation*}
  e(W_h)=r_h-(d-1)h^2+O_d(h^4),
  \qquad
  \Phi_H(p_h,r_h)
  =J_{p_h}(r_h)-\frac{d^3}{3}h^4+O_d(h^6).
\end{equation*}
\end{theorem}

\subsection{The rank-one stationary family near the exceptional endpoint}
\label{sec:rank-one-stationary-family}

We now construct an analytic family of rank-one bipodal
candidates.  The numerical exploration in \Cref{fig:bipodal-parameter-maps}
suggests that rank-one graphons describe the optimizers along a curve of
parameter pairs $(p,r)$ approaching $P_*$.  Motivated by this observation,
we seek rank-one solutions of the full graphon stationarity equations,
together with the block-proportion stationarity condition.
At this stage, the resulting family consists only of stationary
rank-one bipodal candidates; the subsequent arguments establish their optimality and
uniqueness along the corresponding parameter curve.

Throughout this subsection, graphons are assumed to take values in a compact
subinterval of $(0,1)$.  For a rank-one graphon $W=f\otimes f$ with
measurable $f:[0,1]\to[0,1]$, this requires that, for some $\eta>0$,
\[
  \eta\le f\le1-\eta
  \qquad\text{a.e.}
\]
This interiority condition is used only to construct the candidate family in
this subsection.  It ensures that $W+\varepsilon U$ remains a graphon for every
bounded symmetric function $U$ and all sufficiently small $|\varepsilon|$, so
that stationarity can be tested in every bounded direction.  The subsequent
optimality and uniqueness arguments impose no such restriction on competing
graphons.  We say that $W$ is \emph{stationary} for
$I_p-\mu t(H,\,\cdot\,)$ if
\begin{equation}\label{eq:graphon-stationarity}
  \left.\frac{d}{d\varepsilon}\right|_{\varepsilon=0}
  \left[
    I_p(W+\varepsilon U)
    -\mu t(H,W+\varepsilon U)
  \right]=0
\end{equation}
for every such $U$.
For a rank-one graphon $W=f\otimes f$, we say that $W$ \emph{satisfies the
KKT conditions} at $(p,r)$ with multiplier $\mu$, if \eqref{eq:graphon-stationarity}
holds with $\mu\ge0$ and $t(H,W)=r^m$.  This is the graphon
analog of the usual first-order Karush--Kuhn--Tucker condition for
constrained finite-dimensional optimization.

To formulate the separate first-order condition for block proportion,
for $\alpha\in(0,1)$ and $s,t\in(0,1)$, set
\[
  f_{\alpha,s,t}
  :=s\1_{[0,\alpha]}+t\1_{(\alpha,1]},
  \qquad
  W_{\alpha,s,t}:=f_{\alpha,s,t}\otimes f_{\alpha,s,t}.
\]
Changing $\alpha$ moves the boundary between the two vertex classes.  Unlike
the perturbations $W_{\alpha,s,t}+\varepsilon U$ used in graphon stationarity,
replacing $\alpha$ by $\alpha+\varepsilon$, when $s\ne t$, changes the graphon
by an amount bounded away from zero on a set whose measure tends to zero.
Hence this variation is not small in $L^\infty$ and is not covered by graphon
stationarity.  We therefore impose the separate condition
\begin{equation}\label{eq:block-stationarity}
  \left.\frac{d}{d\beta}\right|_{\beta=\alpha}
  \left[
    I_p(W_{\beta,s,t})-\mu t(H,W_{\beta,s,t})
  \right]=0,
\end{equation}
obtained by varying the block proportion while keeping $s$ and $t$ fixed.

The following lemma makes a structural reduction: it shows that bipodality is
forced by the KKT conditions within the interior rank-one class.  The
construction lemma later begins by seeking a bipodal family; this structural
result justifies that choice and removes the need to assume bipodality in
the accompanying local uniqueness assertion.  The proof reduces the stationarity condition
\eqref{eq:graphon-stationarity} to a scalar equation and bounds the
number of its roots.  For this reduction, define the \emph{$d$-th moment} of
$f$ by
\[
  q(f):=\int_0^1 f(x)^d\,dx.
\]
Note that any rank-one graphon $W=f\otimes f$ satisfies
$t(H,W)=q(f)^v$ depending only on $H$ through $d$ and $v$:
\[
  t(H,f\otimes f)
  =\int_{[0,1]^v}\prod_{(i,j)\in E(H)}f(x_i)f(x_j)\,dx_1\cdots dx_v
  =\int_{[0,1]^v}\prod_{i=1}^v f(x_i)^{d}\,dx_1\cdots dx_v
  =q(f)^v.
\]

\begin{lemma}[Bipodality of stationary rank-one graphons]
\label{lem:stationary-rank-one-bipodality}
Let $p,r\in(0,1)$, let $\mu\ge0$, and let
$f:[0,1]\to[0,1]$ be measurable.  Suppose that, for some $\eta>0$,
\[
  \eta\le f\le1-\eta
  \qquad\text{a.e.},
\]
and that $f\otimes f$ satisfies the KKT conditions at $(p,r)$ with multiplier
$\mu$.  If $f$ is not a.e.\ constant, then there exist $\alpha\in(0,1)$ and
$0<s<t<1$ such that, after a measure-preserving relabeling,
\[
  f=f_{\alpha,s,t}
  \qquad\text{a.e.}
\]
\end{lemma}

\begin{proof}
Set $\gamma:=\mu m q(f)^{v-2}$ and $F_{p,\gamma}(z):=J_p'(z)-\gamma z^{d-1}$.
We first claim that
\begin{equation}
\label{eq:rank-one-kkt}
  F_{p,\gamma}(f(x)f(y))=0
  \quad\text{for a.e. }(x,y)\in[0,1]^2.
\end{equation}
Let $U$ be any bounded symmetric function on $[0,1]^2$ and set
$W_\varepsilon=f\otimes f+\varepsilon U$.  The variation of $I_p$ is
\[
  \left.\frac{d}{d\varepsilon}I_p(W_\varepsilon)\right|_{\varepsilon=0}
  =\iint J_p'(f(x)f(y))U(x,y)\,dx\,dy.
\]
Since $H$ is $d$-regular, differentiating one of the $m$ edge factors in
$t(H,W_\varepsilon)$ gives
\[
  \left.\frac{d}{d\varepsilon}t(H,W_\varepsilon)\right|_{\varepsilon=0}
  =\iint m q(f)^{v-2}f(x)^{d-1}f(y)^{d-1}
  U(x,y)\,dx\,dy.
\]
Stationarity for every $U$ therefore yields
\[
  J_p'(f(x)f(y))
  =\mu m q(f)^{v-2}f(x)^{d-1}f(y)^{d-1}
  =\gamma\bigl(f(x)f(y)\bigr)^{d-1}
\]
for almost every $(x,y)$, proving
\eqref{eq:rank-one-kkt}.

We next show that $F_{p,\gamma}(z)=0$ has at most three
solutions for $z\in(0,1)$.  Since $J_p''(z)=1/(z(1-z))$,
\[
  F_{p,\gamma}'(z)
  =\frac1{z(1-z)}-\gamma(d-1)z^{d-2}.
\]
If $\gamma\le0$, then $F_{p,\gamma}'>0$ on $(0,1)$, so
$F_{p,\gamma}$ has at most one zero.  If $\gamma>0$, the critical-point
equation is
\[
  z^{d-1}(1-z)=\frac1{(d-1)\gamma}.
\]
The function $z\mapsto z^{d-1}(1-z)$ is strictly increasing on
$(0,(d-1)/d)$ and strictly decreasing on $((d-1)/d,1)$.  Hence
$F_{p,\gamma}'$ has at most two distinct zeros,
and by Rolle's theorem, $F_{p,\gamma}$ has at most three distinct zeros.

By Fubini's theorem, one can choose $y_0\in[0,1]$ such that
\[
  F_{p,\gamma}(f(x)f(y_0))=0
  \qquad\text{for a.e. }x.
\]
The lower bound on $f$ gives $f(y_0)>0$.  Hence multiplication by
$f(y_0)$ is injective, and the three-root bound for $F_{p,\gamma}$ implies
that $f$ takes at most three values up to null sets.

Suppose that it takes three distinct values $0<a<b<c<1$ up to null sets.
Since \eqref{eq:rank-one-kkt} holds outside a null subset
of $[0,1]^2$, each of
\[
  a^2,\qquad ab,\qquad ac,\qquad bc,\qquad c^2
\]
is a zero of $F_{p,\gamma}$.  They satisfy
$a^2<ab<ac<bc<c^2$, contradicting the three-root bound.  Hence $f$ takes at
most two values.
Therefore, if $f$ is not a.e. constant, it has the form
\[
  f=s\1_A+t\1_{A^c}
\]
for some $0<s<t<1$ and measurable $A$ with $0<\lambda(A)<1$.  Setting
$\alpha:=\lambda(A)$ and relabeling $A$ as $[0,\alpha]$ gives
$f=f_{\alpha,s,t}$ almost everywhere.
\end{proof}

\Cref{lem:stationary-rank-one-bipodality} reduces the problem to a finite-dimensional
one: after a measure-preserving relabeling, the factor of every nonconstant
rank-one graphon satisfying the KKT conditions has the form
$f_{\alpha,s,t}$ for some $\alpha\in(0,1)$ and $0<s<t<1$.  It does not,
however, construct such graphons near $P_*$ or determine their parameters.
The next lemma constructs this analytic family.
Reparametrize $s$ and $t$ by
\[
  u:=\frac{s+t}{2},
  \qquad
  h:=\frac{t-s}{2},
  \qquad
  \gamma:=\mu m q(f_{\alpha,s,t})^{v-2}.
\]
Thus, $u$ is the midpoint of the two factor values, $|h|$ is half their
separation, and $\gamma$ is the coefficient appearing in the scalar
stationarity equation \eqref{eq:rank-one-kkt}.
The three stationarity equations, evaluated at the three graphon values,
determine $u,p,\gamma$; the block-proportion equation then determines
$\alpha$, and the defining relation for $\gamma$ determines the multiplier $\mu$.
Changing $h$ to $-h$ exchanges the two factor values and their proportions,
while remaining quantities are unchanged.

We denote the value of the coefficient $\gamma$ at the singular endpoint by
\[
  \gamma_*:=\left(\frac{d}{d-1}\right)^d=r_*^{-d}.
\]

\begin{lemma}[Analytic rank-one KKT family]
\label{lem:rank-one-kkt-family}
There exist $h_0>0$ and real-analytic functions
\[
  h\longmapsto(u_h,p_h,\gamma_h,\alpha_h)
  \qquad (|h|<h_0)
\]
such that
\[
  (u_0,p_0,\gamma_0,\alpha_0)
  =\left(u_*,p_*,\gamma_*,\frac12\right),
\]
and the following properties hold.  First, for all $|h|<h_0$,
\[
  u_{-h}=u_h,\qquad p_{-h}=p_h,\qquad
  \gamma_{-h}=\gamma_h,\qquad \alpha_{-h}=1-\alpha_h.
\]
Next, for $|h|<h_0$, define the two-valued factor and its rank-one graphon by
\[
  s_h:=u_h-h,
  \qquad 
  t_h:=u_h+h,
  \qquad
  f_h:=f_{\alpha_h,s_h,t_h}, 
  \qquad
  W_h:=f_h\otimes f_h,
\]
and define the $d$-th moment of $f_h$, the associated target density, and the
Lagrange multiplier by
\[
  q_h:=q(f_h)=\alpha_hs_h^d+(1-\alpha_h)t_h^d,
  \qquad
  r_h:=q_h^{2/d}, 
  \qquad
  \mu_h:=\frac{\gamma_h}{m q_h^{v-2}}.
\]
The resulting functions $q_h,r_h,\mu_h$ are real analytic, and for all
$|h|<h_0$, these quantities satisfy
\[
  0<p_h<r_h<1,
  \qquad s_h,t_h\in(0,1),
  \qquad \alpha_h\in(0,1),
  \qquad \mu_h>0,
  \qquad\text{and}\qquad
  t(H,W_h)=r_h^m.
\]
Finally, for $0<|h|<h_0$, the nonconstant rank-one graphon $W_h$ satisfies
the KKT conditions at $(p_h,r_h)$ with multiplier $\mu_h$, and
\eqref{eq:block-stationarity} holds.  Conversely, in the precise
sense stated below, let $W=f\otimes f$ be a nearby nonconstant rank-one
graphon whose factor $f$ is bounded away from $0$ and $1$.  If $W$ satisfies
the KKT conditions with multiplier $\mu$ and
\eqref{eq:block-stationarity} holds, then $W$ agrees, up to
measure-preserving relabeling, with a unique member $W_h$ indexed by
$h\in(0,h_0)$, and $\mu=\mu_h$.
\end{lemma}

The graphons $W_h$ supplied by \Cref{lem:rank-one-kkt-family},
with their measure-preserving relabelings, form the family used in
\Cref{thm:singular-endpoint}.  Before proving the lemma, we clarify the
family at the singular endpoint and the local uniqueness assertion.
At $h=0$, the definitions in the lemma give
\[
  s_0=t_0=u_*,
  \qquad
  q_0=u_*^d,
  \qquad
  r_0=u_*^2=r_*,
  \qquad
  \mu_0=\frac1{m r_*^m},
  \qquad
  W_0\equiv r_*.
\]
Thus, $(p_0,r_0)=P_*$, so the family passes through the constant graphon
$W_0\equiv r_*$ at the singular endpoint.  Also, the two factor values $s_0$ and
$t_0$ coincide, as do the three values $s_0^2,s_0t_0,t_0^2$ taken by the
graphon $W_0$.
For the local uniqueness assertion, consider a graphon $W=f\otimes f$ and a
multiplier $\mu$ appearing in the converse assertion of the lemma.  By
\Cref{lem:stationary-rank-one-bipodality}, after a measure-preserving relabeling its
factor has the form $f_{\alpha,s,t}$ with
$\alpha\in(0,1)$ and $0<s<t<1$.  If
$(s,t,p,r,\mu,\alpha)$ is sufficiently close to
\[
  \left(u_*,u_*,p_*,r_*,\frac1{m r_*^m},\frac12\right),
\]
then $h:=(t-s)/2$ belongs to $(0,h_0)$ and
\[
  (s,t,p,r,\mu,\alpha)
  =(s_h,t_h,p_h,r_h,\mu_h,\alpha_h).
\]

\begin{proof}[Proof outline]
Write
\[
  \ell(z):=\log\frac{1-z}{z},
  \qquad
  \ell_*:=\ell(p_*)=\frac{d}{d-1}-\log(d-1).
\]
By \Cref{lem:stationary-rank-one-bipodality}, every nonconstant rank-one graphon
satisfying the KKT conditions has, after a measure-preserving relabeling, a
two-valued factor.  Motivated by this reduction, we seek a family of the form
\[
  s:=u-h,
  \qquad
  t:=u+h,
  \qquad
  f:=f_{\alpha,s,t},
  \qquad
  W:=f\otimes f.
\]
Since $f\otimes f$ takes the three values $s^2,st,t^2$, the scalar
stationarity equation \eqref{eq:rank-one-kkt} requires
\begin{equation}\label{eq:three-value-kkt}
  \ell(p)-\ell(s^2)=\gamma s^{2d-2},\qquad
  \ell(p)-\ell(st)=\gamma(st)^{d-1},\qquad
  \ell(p)-\ell(t^2)=\gamma t^{2d-2}.
\end{equation}
We therefore begin by solving this necessary system for $u,p,\gamma$ as
functions of $h$; the block proportion $\alpha$ does not occur and will be
determined later.
At $h=0$, these equations have the solution $(u,p,\gamma)=(u_*,p_*,\gamma_*)$.
For $h\ne0$, subtracting consecutive equations in \eqref{eq:three-value-kkt} eliminates
$\ell(p)$ and expresses $\gamma$ as each of the two divided slopes across
$s^2,st,t^2$.  Although each slope appears as $0/0$ at $h=0$, both extend
analytically.  Let $A(u,h)$ denote the difference of the two divided slopes,
so the system requires $A(u,h)=0$.  The symmetry $h\mapsto-h$ makes $A$ odd in $h$,
and hence $A(u,h)=hB(u,h)$ for some analytic $B$.  A direct calculation shows that
$B(u,0)$ has the unique nondegenerate zero $u=u_*$.  The analytic implicit
function theorem therefore gives $h_0>0$ and a unique even analytic solution
$u=u_h$ for $|h|<h_0$ satisfying $B(u_h,h)=A(u_h,h)=0$.  The common divided slope then defines the even function
$\gamma=\gamma_h$, and the middle equation in
\eqref{eq:three-value-kkt} determines the even function $p=p_h$.

We now choose the block proportion $\alpha$ as a function of $h$.
Once $\alpha$ is fixed, set
\[
  q:=\alpha s_h^d+(1-\alpha)t_h^d,
  \qquad
  r:=q^{2/d},
  \qquad
  \mu:=\frac{\gamma_h}{m q^{v-2}},
  \qquad
  f:=f_{\alpha,s_h,t_h},
  \qquad
  W:=f\otimes f.
\]
Also define the function
\[
  \mathcal M_h(z):=J_{p_h}(z)-\frac{\gamma_h}{d}z^d
  \qquad (0 < z < 1).
\]
Evaluating the derivative in \eqref{eq:block-stationarity} and
solving the resulting linear equation gives
\begin{equation}\label{eq:block-proportion-formula}
  \alpha_h:=
  \frac{\mathcal M_h(t_h^2)-\mathcal M_h(s_ht_h)}
  {\left\{\mathcal M_h(t_h^2)-\mathcal M_h(s_ht_h)\right\}+
  \left\{\mathcal M_h(s_h^2)-\mathcal M_h(s_ht_h)\right\}}.
\end{equation}
To analyze this quotient as $h\to0$, we factor $\mathcal M_h'$.  By
\eqref{eq:three-value-kkt}, the three values
$s_h^2,s_ht_h,t_h^2$ are zeros of $\mathcal M_h'$, and
\begin{equation}\label{eq:mh-derivative-factorization}
  \mathcal M_h'(z)
  =(z-s_h^2)(z-s_ht_h)(z-t_h^2)R(h,z),
\end{equation}
where $R$ is jointly analytic near $(0,r_*)$.
At $h=0$, the three zeros coincide at $r_*$, and direct calculation gives
\[
  \mathcal M_0'(z)
  =k_3(z-r_*)^3+O_d((z-r_*)^4),
  \qquad
  k_3:=\frac{d^5}{6(d-1)^2}>0.
\]
Thus $R(0,r_*)=k_3$.  Since every $z\in[s_h^2,t_h^2]$ satisfies
$|z-r_*|=O_d(|h|)$, joint analyticity of $R$ gives
\[
  R(h,z)=k_3+O_d(|z-r_*|+|h|)=k_3+O_d(|h|).
\]
Substituting this into \eqref{eq:mh-derivative-factorization}
and integrating over $[s_h^2,s_ht_h]$ and $[s_ht_h,t_h^2]$,
and using the expansions of $s_h$ and $t_h$ in powers of $h$,
show that the two differences in \eqref{eq:block-proportion-formula} vanish to order exactly four
and have the same leading term.  Canceling their common factor $h^4$ extends
the quotient analytically to $h=0$, where it equals $1/2$.  Finally,
$h\mapsto-h$ exchanges the two differences, so
$\alpha_{-h}=1-\alpha_h$.  The resulting analytic function $\alpha_h$ then
determines the analytic functions $q_h,r_h,\mu_h$, and the rank-one
factors $f_h$ and the bipodal graphons $W_h$.

The definitions of $q_h,r_h,\mu_h$ now give
$t(H,W_h)=q_h^v=r_h^m$.  The three equations
\eqref{eq:three-value-kkt} verify the stationarity condition
\eqref{eq:graphon-stationarity}, while
\eqref{eq:block-proportion-formula} verifies
\eqref{eq:block-stationarity}.
After decreasing $h_0$, continuity gives $0<p_h<r_h<1$,
$s_h,t_h,\alpha_h\in(0,1)$, and $\gamma_h,q_h>0$, hence $\mu_h>0$.
Thus, $(p_h,r_h)$ is admissible and $W_h$ is an interior bipodal graphon.

Finally, \Cref{lem:stationary-rank-one-bipodality} reduces any nearby rank-one graphon family in the
converse assertion to the same bipodal form that we selected in the beginning of the proof.
Repeating the construction, uniqueness in the implicit function theorem and in the linear
block-proportion equation forces the graphon, up to relabeling, to be
$W_h$ for a unique $h>0$, with multiplier $\mu_h$.
Full details are given in
\Cref{app:rank-one-kkt-family}.
\end{proof}

\begin{remark}[Dependence on $d$ and $H$]\label{rmk:rank-one-family-universality}
The equations \eqref{eq:three-value-kkt} and \eqref{eq:block-proportion-formula} depend only on $d$.
Thus, for fixed $d$, the local parameter curve $(p_h,r_h)$ and graphon
family $W_h$, up to relabeling, are independent of the particular
$d$-regular graph $H$.  The multiplier $\mu_h=\gamma_h/(m q_h^{v-2})$
and the interval on which we will establish global optimality and uniqueness
in the subsequent sections may depend on $H$.
\end{remark}

\subsection{Comparison with the constant graphon}
\label{sec:constant-graphon-comparison}

The constant graphon $W\equiv r_h$ has $H$-density $r_h^m$ and is therefore
feasible for the same constraint as $W_h$, where $W_h$ is the member of the
analytic rank-one KKT family constructed in
\Cref{lem:rank-one-kkt-family}.  We will prove that
\[
  I_{p_h}(W_h)<I_{p_h}(W\equiv r_h)=J_{p_h}(r_h),
\]
and compute the asymptotic expansion of the
relative entropy difference for all sufficiently small $h>0$.  At $h=0$, the
two graphons coincide, so their $I_{p_h}$-values are equal.  The symmetries in
\Cref{lem:rank-one-kkt-family} make the relative entropy difference
an even analytic function of $h$, and the calculation below shows that its
leading term is of order $h^4$.  The sign of the difference for small $h>0$ is therefore
determined by the coefficient of $h^4$.  To compute that coefficient, we
first record the expansions of the parameters that will be used below.
Write
\[
  \ell(z):=\log\frac{1-z}{z}, \qquad
  \ell_*:=\ell(p_*)=\frac d{d-1}-\log(d-1), \qquad
  \Lambda_h:=\ell(p_h)-\ell_*.
\]
Here, $\Lambda_h$ represents the slope of the affine difference
\begin{equation}\label{eq:entropy-parameter-shift}
  J_{p_h}(z)-J_{p_*}(z)=\Lambda_h z+J_{p_h}(0)-J_{p_*}(0).
\end{equation}

\begin{lemma}[Expansions along the rank-one KKT family]
\label{lem:rank-one-parameter-expansions}
The analytic family constructed in
\Cref{lem:rank-one-kkt-family} satisfies
\begin{equation}\label{eq:rank-one-parameter-expansions}
\begin{aligned}
  u_h&=u_*+\frac{5-3d}{6u_*}h^2+O_d(h^4),
  &&&
  \Lambda_h&=\frac{2d^3}{3(d-1)}h^2+O_d(h^4),\\
  \alpha_h&=\frac12+\frac{2d^2u_*}{3(d-1)}h+O_d(h^3),
  &&&
  r_h&=r_*-\frac{2(4d-1)}3h^2+O_d(h^4).
\end{aligned}
\end{equation}
\end{lemma}

These expansions immediately locate the family relative to
$P_*=(p_*,r_*)$.  For all sufficiently small $h>0$, the expansion of $r_h$
gives $r_h<r_*$.  Moreover, since $\ell$ is strictly decreasing,
$\Lambda_h>0$ implies $p_h<p_*$.

\begin{proof}[Proof outline]
The symmetries in \Cref{lem:rank-one-kkt-family} show that
$u_h,\gamma_h$, and $\Lambda_h$ have expansions in even powers of $h$, while
$\alpha_h-1/2$ has an expansion in odd powers.  Let $U_2, G_2, L_2$
denote the quadratic coefficients of $u_h,\gamma_h$, and $\Lambda_h$,
respectively, and substitute the expansions of $s_h^2,s_ht_h,t_h^2$ into
\eqref{eq:three-value-kkt}.
The middle equation at order $h^2$ relates $G_2$ and $L_2$,
and either outer equation at order $h^3$ determines $G_2$,
and the comparison at order $h^4$ relates $U_2$ and $G_2$.
Solving these three equations gives the first two expansions in
\eqref{eq:rank-one-parameter-expansions}.

Now we determine the expansions of $\alpha_h$ and $r_h$.
Since $R(h,z)$ is even in $h$, its Taylor expansion gives
\[
  R(h,z)=k_3+k_4(z-r_*)+O_d(h^2)
  \quad (z\in[s_h^2,t_h^2]),
  \qquad
  k_4:=\frac{5d^6(d-2)}{24(d-1)^3}.
\]
Using this expansion and integrating the factorization
\eqref{eq:mh-derivative-factorization} over $[s_h^2,s_ht_h]$ and
$[s_ht_h,t_h^2]$ gives the two differences in \eqref{eq:block-proportion-formula}
through order $h^5$.  Their asymmetry at order $h^5$ determines the linear
term of $\alpha_h$.  The expansions of $q_h$ and $r_h$ then follow from
their definitions.
Full details are given in \Cref{app:rank-one-parameter-expansions}.
\end{proof}

We now turn to the comparison stated at the beginning of this subsection.

\begin{lemma}[Comparison with the constant graphon]
\label{lem:constant-graphon-comparison}
For sufficiently small $h>0$,
\[
  e(W_h)=r_h-(d-1)h^2+O_d(h^4),
  \qquad
  I_{p_h}(W_h)=J_{p_h}(r_h)-\frac{d^3}{3}h^4+O_d(h^6).
\]
In particular, the constant graphon
$W\equiv r_h$ is feasible but has strictly larger $I_{p_h}$-value,
and consequently $p_h<\pc(r_h)$ for all sufficiently small $h>0$.
\end{lemma}

The affine identity \eqref{eq:entropy-parameter-shift} separates the change in $p$
from the comparison at $p_*$:
\[
  I_{p_h}(W_h)-J_{p_h}(r_h)
  =\{I_{p_*}(W_h)-J_{p_*}(r_h)\}+\Lambda_h\{e(W_h)-r_h\}.
\]
The second term is evaluated using \Cref{lem:rank-one-parameter-expansions}.
For the first difference, we return to the supporting-line geometry of
\Cref{thm:scalar-lz-boundary}.  Its proof obtains the quadratic separation
in \eqref{eq:contact-quadratic-separation} by Taylor-expanding the difference between
$\varphi_{\pc(r),d}(x)=J_{\pc(r)}(x^{1/d})$ and its tangent line $\ell_r(x)$
near the two contact points.  We use the same approach at the exceptional
endpoint, where these contacts coincide.

By the continuous extensions in \Cref{thm:scalar-lz-boundary},
$\pc(r)\to p_*$ and $\sm(r)\to r_*$ as $r\to r_*$.  The tangent lines
$\ell_r$ therefore extend continuously to the tangent of
$\varphi_{p_*,d}$ at $x=r_*^d$:
\[
  \ell_{r_*}(x)=J_{p_*}(r_*)+\beta_d(x-r_*^d).
\]
Here $\beta_d$ denotes the slope of this tangent, which the chain rule gives as
\[
  \beta_d:=\varphi_{p_*,d}'(r_*^d)
  =\frac{J_{p_*}'(r_*)}{d r_*^{d-1}}
  =\frac{\gamma_*}{d}
  =\frac{d^{d-1}}{(d-1)^d}.
\]
Following \eqref{eq:supporting-cost-gap}, define the \emph{supporting cost gap} by
\begin{equation}\label{eq:endpoint-supporting-gap}
  \Gamma_d(z)
  :=\varphi_{p_*,d}(z^d)-\ell_{r_*}(z^d)
  =J_{p_*}(z)-J_{p_*}(r_*)-\beta_d\{z^d-r_*^d\}
  \qquad(0\le z\le1).
\end{equation}
Thus, $\Gamma_d$ is the continuous extension to $r=r_*$ of the function
$g_r$ used there.  Its Taylor expansion describes the separation from
the tangent at their common contact point; as shown below, the leading term is quartic.

By \eqref{eq:endpoint-supporting-gap}, we can express the cost difference
at $p_*$ in terms of $\Gamma_d$ as
\[
  I_{p_*}(W_h)-J_{p_*}(r_h)
  =\int\Gamma_d(W_h)-\Gamma_d(r_h)+\beta_d\left(\int W_h^d-r_h^d\right).
\]
Adding the affine correction from $p_*$ to $p_h$ gives
\begin{equation}\label{eq:constant-comparison-decomposition}
  I_{p_h}(W_h)-J_{p_h}(r_h)
  =\int\Gamma_d(W_h)-\Gamma_d(r_h)
  +\beta_d\left(\int W_h^d-r_h^d\right)
  +\Lambda_h\{e(W_h)-r_h\}.
\end{equation}
Since $\int W_h^d=(\int f_h^d)^2=q_h^2=r_h^d$, the $d$-th moment term vanishes, and
\begin{equation}\label{eq:constant-comparison-identity}
  I_{p_h}(W_h)-J_{p_h}(r_h)
  =\int\Gamma_d(W_h)-\Gamma_d(r_h) +\Lambda_h\{e(W_h)-r_h\}.
\end{equation}
The proof first evaluates $\Gamma_d$ by its quartic expansion,
then expands all the parameters using \Cref{lem:rank-one-parameter-expansions}.

\begin{proof}
Tangency gives $\Gamma_d(r_*)=\Gamma_d'(r_*)=0$.  At the exceptional
endpoint, direct differentiation further gives
$\Gamma_d''(r_*)=\Gamma_d'''(r_*)=0$ and
$\Gamma_d^{(4)}(r_*)=d^5/(d-1)^2>0$.
Taylor's theorem therefore yields, as $z\to r_*$,
\begin{equation}\label{eq:endpoint-gap-expansion}
  \Gamma_d(z)
  =\frac{d^5}{24(d-1)^2}(z-r_*)^4+O_d(|z-r_*|^5).
\end{equation}
By the parameter expansions in \eqref{eq:rank-one-parameter-expansions},
\[
  s_h^2-r_*=-2u_*h+O_d(h^2),
  \qquad
  t_h^2-r_*=2u_*h+O_d(h^2),
  \qquad
  s_ht_h-r_*=O_d(h^2),
  \qquad
  r_h-r_*=O_d(h^2).
\]
Each of the four rectangles determined by $A_h$ and $A_h^c$ in
$[0,1]^2$ has area $1/4+O_d(h)$.
Thus the quartic expansion and $u_*^2=(d-1)/d$ yield
\begin{equation}\label{eq:family-gap-expansion}
  \int\Gamma_d(W_h)=\frac{d^3}{3}h^4+O_d(h^6),
  \qquad
  \Gamma_d(r_h)=O_d(h^8).
\end{equation}
The first remainder improves from $O_d(h^5)$ to $O_d(h^6)$
because $\int\Gamma_d(W_h)$ is an even analytic function of $h$.
Finally, $e(W_h)=\{u_h+h(1-2\alpha_h)\}^2$ and the parameter expansions give
\[
  e(W_h)-r_h=-(d-1)h^2+O_d(h^4),
  \qquad
  \Lambda_h\{e(W_h)-r_h\}=-\frac{2d^3}{3}h^4+O_d(h^6).
\]
Substituting these expansions into \eqref{eq:constant-comparison-identity}
proves the lemma.
\end{proof}



\subsection{Localization and rank-one reduction}
\label{sec:localization-rank-one}

We now take any feasible graphon $W$ satisfying
$I_{p_h}(W)\le I_{p_h}(W_h)$ and prepare it for comparison with $W_h$.
This has two stages.  First, the $I_{p_h}$-comparison forces $W$ to be close
in $L^4$ to the constant graphon $W\equiv r_*$; this is the localization
step.  Second, we choose a rank-one part $f\otimes f$ and write
$W=f\otimes f+E$.  The first-stage $L^4$-localization then shows that the factor
$f$ is close to $u_*$ and that the error $E$ is small.  The lower bound for
the auxiliary Lagrangian difference obtained in
\Cref{sec:auxiliary-lagrangian} will control the distributions of
the values of $f$, while the arbitrary-graphon comparison in
\Cref{sec:graphon-comparison} will control $E$.
Thus, the decomposition is the form in which the localization is used in the
remainder of the proof.

For the first stage, we use the supporting cost gap $\Gamma_d$ from
\eqref{eq:endpoint-supporting-gap}.  The mechanism is to bound
$\int\Gamma_d(W)$ using the $I_{p_h}$-comparison and then convert that bound
into an $L^4$-bound on $W-r_*$.  Its quartic expansion gives the required
control near $r_*$; to extend this control to all of $[0,1]$, we also need
strict positivity away from $r_*$.  This follows from the supporting-line
interpretation above: by \Cref{lem:convexity-defect}, $\varphi_{p_*,d}$ is
convex on $[0,1]$ and its tangent $\ell_{r_*}$ has no other contact.
Therefore,
\begin{equation}\label{eq:endpoint-gap-positivity}
  \Gamma_d(z)\ge0,
  \qquad
  \Gamma_d(z)=0\iff z=r_*.
\end{equation}
This makes $\int\Gamma_d(W)$ a nonnegative measure of how far $W$ is from
the constant graphon $W\equiv r_*$.  For $L^4$-localization, we use the
stronger pointwise bound
\begin{equation}\label{eq:endpoint-gap-quartic-bound}
  \Gamma_d(z)\ge c_d|z-r_*|^4
  \qquad(0\le z\le1)
\end{equation}
for some $c_d>0$.  Indeed, the ratio $\Gamma_d(z)/|z-r_*|^4$ has a positive
limit at $r_*$ by \eqref{eq:endpoint-gap-expansion}, hence a
positive lower bound in a neighborhood of $r_*$.  On the compact complement
of this neighborhood, it has a positive minimum by continuity and
\eqref{eq:endpoint-gap-positivity}.  This proves
\eqref{eq:endpoint-gap-quartic-bound}, whose integral form is
\[
  \int\Gamma_d(W)\ge c_d\|W-r_*\|_4^4.
\]
Thus, $\int\Gamma_d(W)=O_d(h^4)$ implies the desired localization
$\|W-r_*\|_4=O_d(h)$.

To derive this integral bound from the $I_{p_h}$-comparison, we relate
$I_{p_h}$ to the support cost gap $\Gamma_d$.  Repeating the calculation leading to
\eqref{eq:constant-comparison-decomposition}, instead using $W$ and $W_h$, gives
\begin{equation}\label{eq:graphon-cost-decomposition}
  I_{p_h}(W)-I_{p_h}(W_h)
  =\int\Gamma_d(W)-\int\Gamma_d(W_h)
  +\beta_d\left(\int W^d-r_h^d\right)
  +\Lambda_h\{e(W)-e(W_h)\}.
\end{equation}
The proof of the next lemma uses this decomposition to obtain the required
bound on $\int\Gamma_d(W)$.
Briefly, the $I_{p_h}$-comparison makes the right-hand side nonpositive.
Using \eqref{eq:endpoint-gap-quartic-bound}, we absorb the edge-density
correction into half of $\int\Gamma_d(W)$ and an $O_d(h^4)$ remainder.
The remaining quantities $\int\Gamma_d(W)$ and $\int W^d-r_h^d$ are
nonnegative, the latter by feasibility.  Since $\int\Gamma_d(W_h)=O_d(h^4)$,
both must be $O_d(h^4)$, yielding the desired localization.

The preceding discussion gives the first stage of localization, but this
alone does not suffice for the later comparison with $W_h$, since $W$ need
not be rank-one.  The lemma below therefore extracts a
rank-one part $f\otimes f$ and writes the remainder as $E$.  For a symmetric
kernel $K$ and a function $g$, write
\[
  (T_Kg)(x):=\int_0^1K(x,y)g(y)\,dy.
\]
The factor is chosen so that $T_E(f^{d-1})=0$.  This identity eliminates the
terms linear in $E$ from the later moment and homomorphism-density expansions.
The following lemma records both the localization and the rank-one reduction.

\begin{lemma}[Localization and rank-one reduction]
\label{lem:localization-rank-one}
There exist $C_d,C_{d,m}>1$ and $h_0>0$ such that, for every
$0<h<h_0$, every graphon $W$ satisfying $t(H,W)\ge r_h^m$ and
$I_{p_h}(W)\le I_{p_h}(W_h)$ satisfies
\begin{equation}\label{eq:graphon-quartic-localization}
  \|W-r_*\|_4\le C_dh.
\end{equation}
Moreover, $W$ admits a decomposition $W=f\otimes f+E$ for some
$f\in L^d([0,1])$ and $E\in L^2([0,1]^2)$ satisfying
\begin{equation}\label{eq:rank-one-orthogonality}
  T_W(f^{d-1})=q(f)f,
  \qquad
  T_E(f^{d-1})=0,
  \qquad
  q(f)>0,
  \qquad
  0\le f(x) \le 2
  \quad\text{for a.e. }x.
\end{equation}
Also, the function $f$ is close to the constant function $u_*$, the error
$E$ is small, the $d$-th moment $q(f)$ is close to $q_h$, and the
homomorphism-density error is bounded by a cubic term in $\|E\|_2$:
\begin{equation}\label{eq:rank-one-reduction-bounds}
  \begin{gathered}
  \|f-u_*\|_4\le C_dh,
  \qquad
  \|E\|_2 \le C_dh,
  \qquad
  \left|q(f)-q_h\right|\le C_dh^2,
  \qquad
  |t(H,W)-q(f)^v| \le C_{d,m} \|E\|_2^3.
  \end{gathered}
\end{equation}
\end{lemma}

\begin{proof}
We first localize $W$ in $L^4$ around the constant graphon $W\equiv r_*$,
then construct its rank-one part $f\otimes f$ and an error $E$ satisfying
$T_E(f^{d-1})=0$, and finally bound the moment and homomorphism-density
errors.

\medskip
\noindent\emph{Localization of the graphon.}
Combining the $I_{p_h}$-comparison with \eqref{eq:graphon-cost-decomposition} gives
\[
  0\ge I_{p_h}(W)-I_{p_h}(W_h)
  =\int\Gamma_d(W)-\int\Gamma_d(W_h)
  +\beta_d\left(\int W^d-r_h^d\right)
  +\Lambda_h\{e(W)-e(W_h)\}.
\]
Generalized H\"older \eqref{eq:generalized-holder} and feasibility give $\int W^d\ge r_h^d$.
To absorb the edge-density correction, convexity, Taylor's theorem, and
\eqref{eq:endpoint-gap-quartic-bound} give
\[
  0\le z^d-r_*^d-d r_*^{d-1}(z-r_*)
  \le C_d|z-r_*|^2
  \le C_d\Gamma_d(z)^{1/2}
  \qquad(0\le z\le1).
\]
Applying the upper bound to $W$ and the lower bound to $W_h$, then using
$\int W_h^d=r_h^d$, yields
\[
  e(W)-e(W_h)
  \ge\frac{\int W^d-r_h^d}{d r_*^{d-1}}
   -C_d\int\Gamma_d(W)^{1/2}
  \ge-C_d\int\Gamma_d(W)^{1/2}.
\]
By \eqref{eq:rank-one-parameter-expansions} and \eqref{eq:family-gap-expansion},
we have $\Lambda_h=O_d(h^2)$ and $\int\Gamma_d(W_h)=O_d(h^4)$.
Substitution into the cost comparison gives
\begin{align*}
  \beta_d \left(\int W^d-r_h^d\right)+\int\Gamma_d(W)
  &\le C_dh^4+C_dh^2\int\Gamma_d(W)^{1/2}\\
  &\le C_dh^4+C_dh^2\left(\int\Gamma_d(W)\right)^{1/2}\\
  &\le \left(C_d+\frac{C_d^2}{2}\right)h^4+\frac12\int\Gamma_d(W).
\end{align*}
Absorbing the last term and using $\beta_d>0$ give
\begin{equation}\label{eq:moment-and-cost-gap-bounds}
  0\le\int W^d-r_h^d\le C_dh^4,
  \qquad
  0\le\int\Gamma_d(W)\le C_dh^4.
\end{equation}
Thus, as discussed above, using \eqref{eq:endpoint-gap-quartic-bound} and $\int\Gamma_d(W)=O_d(h^4)$
gives the desired localization \eqref{eq:graphon-quartic-localization}.

\medskip
\noindent\emph{Construction of the decomposition.}
We first construct $f$ and $E$ satisfying \eqref{eq:rank-one-orthogonality}.
Put $d':=d/(d-1)$ so that $(1/d)+(1/d')=1$.
To see the integral operator $T_W:L^{d'}\longrightarrow L^d$ is compact,
choose finite-step graphons $S_n$ satisfying
$\|W-S_n\|_d\to0$.  For each $n$, choose a finite partition
$A_1^{(n)},\ldots,A_{N_n}^{(n)}$ of $[0,1]$ on whose rectangles $S_n$ is
constant, and write
\[
  S_n(x,y)=\sum_{i,j=1}^{N_n}c_{ij}^{(n)}
  \1_{A_i^{(n)}}(x)\1_{A_j^{(n)}}(y),
  \qquad
  (T_{S_n}g)(x)=\sum_{i=1}^{N_n} \left(\sum_{j=1}^{N_n}c_{ij}^{(n)}\int_{A_j^{(n)}}g\right) \1_{A_i^{(n)}}(x).
\]
Thus, the range of $T_{S_n}$ is contained in the span of the finitely many
functions $\1_{A_1^{(n)}},\ldots,\1_{A_{N_n}^{(n)}}$, so $T_{S_n}$ is compact.
Moreover, H\"older's inequality gives
\[
  \|(T_W-T_{S_n})g\|_d
  \le\|W-S_n\|_d\|g\|_{d'}.
\]
Thus, $S_n\to W$ ensures that $T_{S_n}g\to T_Wg$ uniformly over
$\|g\|_{d'}\le1$.  In other words, $T_{S_n}\to T_W$ in operator norm.
Since the space of compact
operators is closed in the operator norm, $T_W$ is compact.

Now we consider the quadratic form
\[
  \mathcal Q_W(g):= \langle g,T_Wg\rangle =\iint W(x,y)g(x)g(y)\,dx\,dy
  \qquad
  (g \in L^{d'}([0,1]),\, \|g\|_{d'}\le1).
\]
Since $1<d'<\infty$, the closed unit ball of $L^{d'}$ is weakly compact,
and the compactness of $T_W$ makes $\mathcal Q_W$ weakly continuous on this ball.
Thus $\mathcal Q_W$ attains a maximum there at some $g_0$;
let $\theta:=\mathcal Q_W(g_0)$ be the maximum.
The localization bound on $W$ and H\"older's inequality give $\theta \in [r_*/2,1]$,
after decreasing $h_0$, since
\[
  \frac{r_*}{2} \le r_*-\|W-r_*\|_1 \le \int W=\mathcal Q_W(1)\le\theta
  = \mathcal Q_W(g_0)\le\|W\|_d\|g_0\|_{d'}^2\le 1.
\]
Since $W\ge0$, taking the absolute value cannot decrease $\mathcal Q_W$,
and positivity of $\theta$ and homogeneity force every maximizer to have norm
one.  We may therefore choose a maximizer $g_0\ge0$ with $\|g_0\|_{d'}=1$.

Since \(d'>1\), the \(L^{d'}\)-unit sphere is smooth.  Therefore, applying the
Lagrange-multiplier theorem to the maximization of $\mathcal Q_W$ on this sphere
gives \(T_Wg_0=\mu g_0^{d'-1}\) for some constant \(\mu\). This gives
\[
  \theta=\mathcal Q_W(g_0)=\int g_0\,T_Wg_0
  =\mu\int g_0^{d'}=\mu\|g_0\|_{d'}^{d'}=\mu.
\]
Hence $T_Wg_0=\theta g_0^{1/(d-1)}$.
Define $f:=\theta^{1/2}g_0^{1/(d-1)}$ and, for the remainder of the proof,
abbreviate $q:=q(f)=\int f^d=\theta^{d/2}>0$.  Then
\begin{align*}
  T_W(f^{d-1})
  &=T_W(\theta^{(d-1)/2}g_0)
  =\theta^{(d+1)/2}g_0^{1/(d-1)}
  =\theta^{d/2}f=qf,
  \\
  T_{f\otimes f}(f^{d-1})
  &=\left(\int f(y)^d\,dy\right)f
  =qf.
\end{align*}
Let $E:=W-f\otimes f$.  Subtracting these identities gives $T_E(f^{d-1})=0$.
Also, $\theta \in [r_*/2,1]$ implies $q = \theta^{d/2} \in [(r_*/2)^{d/2},1]$.
Consequently, since $f\ge0$,
\[
  qf(x)=\int W(x,y)f(y)^{d-1}\,dy
  \le\int f^{d-1}\le q^{(d-1)/d},
  \qquad
  0\le f(x)\le q^{-1/d}\le\sqrt{\frac{2}{r_*}} \le 2
  \quad\text{for a.e. }x.
\]

\medskip
\noindent\emph{Factor and error bounds.}
Substituting $W=r_*+(W-r_*)$ into the first identity of
\eqref{eq:rank-one-orthogonality} gives
\[
  qf=T_W(f^{d-1})
  =r_*\int f^{d-1}+T_{W-r_*}(f^{d-1}).
\]
The second term on the right-hand side is small, since
$\|T_{W-r_*}(f^{d-1})\|_4 \le\|W-r_*\|_4\|f^{d-1}\|_{4/3} \le C_d h$.
Thus, to show that $f$ is close to $u_*$, it remains to compare the
first term with $qu_*$.  Define $b$ by
\[
  b:=\frac{r_*\int f^{d-1}}q,
  \qquad
  f-b=q^{-1}T_{W-r_*}(f^{d-1}),
  \qquad
  \|f-b\|_4\le C_dh.
\]
The last inequality follows from the bound on the second term above and the
uniform lower bound on $q$.  The uniform bounds on $f$ and $q$ also give a
two-sided uniform bound on $b$:
\[
  \frac{q}{C_d}\le\int f^{d-1}\le q^{(d-1)/d},
  \qquad
  C_d^{-1}\le b\le C_d.
\]
Lipschitz continuity of $x\mapsto x^k$ for $k\in\{d-1,d\}$
on a bounded interval and H\"older's inequality give
\[
  \left|\int f^k-b^k\right|
  \le \|f^k-b^k\|_4
  \le C_d\|f-b\|_4
  \le C_dh
  \qquad(k=d-1,d).
\]
Taking $k=d$ and $k=d-1$, respectively, and using
$q=\int f^d$ and $qb=r_*\int f^{d-1}$ gives
\[
  |q-b^d|\le C_dh,
  \qquad
  |qb-r_*b^{d-1}|\le C_dh.
\]
Since $b$ is uniformly bounded away from zero and $r_*=u_*^2$, these bounds imply
\[
  |b^2-r_*| = \frac{\left| (qb-r_*b^{d-1})-b(q-b^d) \right|}{b^{d-1}} \le C_dh,
  \qquad
  |b-u_*| \le C_dh.
\]
Combining this with the bound on $f-b$ gives
\[
  \|f-u_*\|_4\le \|f-b\|_4+\|b-u_*\|_4\le C_dh.
\]
Finally, $E=W-(f\otimes f)$ is small since $W$ and $f \otimes f$ are both close to $r_*$ in $L^4$:
\begin{equation*}
  \|E\|_2\le\|E\|_4
  \le \|W-r_*\|_4+\|(f\otimes f)-r_*\|_4
  \le \|W-r_*\|_4+(\|f\|_4+u_*)\|f-u_*\|_4
  \le C_dh.
\end{equation*}

\medskip
\noindent\emph{Moment and homomorphism-density errors.}
The binomial expansion of $W^d=(f\otimes f+E)^d$ gives
\begin{equation*}
  \left| \int W^d-q^2 \right|
  = \left| \sum_{k=2}^d\binom dk\int(f\otimes f)^{d-k}E^k \right|
  \le C_d\sum_{k=2}^d\int|E|^k \le C_d\|E\|_2^2.
\end{equation*}
The $k=0$ term cancels with $q^2$, and the $k=1$ term vanishes by orthogonality $T_E(f^{d-1})=0$.
Together with the first bound in \eqref{eq:moment-and-cost-gap-bounds} and $q_h^2=r_h^d$,
\[
  |q-q_h| \le C_d\left|q^2-q_h^2\right|
  \le C_d \left( \left| \int W^d-q^2 \right| + \left| \int W^d - r_h^d \right| \right)
  \le C_dh^2.
\]
Also, expand $t(H,f\otimes f+E)$ according to the edges on which $E$ and $f\otimes f$ are used:
\[
  t(H,W)=\sum_{F\subseteq E(H)}\int
  \prod_{\{i,j\}\in F}E(x_i,x_j)
  \prod_{\{i,j\}\in E(H)\setminus F}f(x_i)f(x_j)\,d\mathbf x.
\]
The $F=\varnothing$ term corresponds to $t(H,f\otimes f)=q^v$, so it remains to bound
the terms with $F\ne\varnothing$.
If a vertex $i$ is incident to exactly one $E$-edge $\{i,j\}$, then integrating
first in $x_i$ gives
\[
  \int E(x_i,x_j)f(x_i)^{d-1}\,dx_i
  =T_E(f^{d-1})(x_j)=0,
\]
so that term vanishes.  Thus, every remaining nonzero term contains a cycle of
$E$-edges, of length $\ell\ge3$ because $H$ is simple.  Since the other
factors are uniformly bounded, we may bound them by a constant and integrate
out all variables outside a chosen cycle.  Thus, the absolute value of the term
is at most
\[
  C_{d,m}\int_{[0,1]^\ell}\prod_{k=1}^\ell|E(x_k,x_{k+1})|\,dx_1\cdots dx_\ell
  \le C_{d,m}\|E\|_2^\ell\le C_{d,m}\|E\|_2^3,
  \qquad x_{\ell+1}:=x_1,
\]
where repeated Cauchy--Schwarz gives the first norm bound, and the last one
uses $\ell\ge3$ and the uniform boundedness of $\|E\|_2$.  Summing the
absolute values of the finitely many terms gives
\begin{equation*}
  |t(H,W)-q^v|\le C_{d,m}\|E\|_2^3.
\end{equation*}
\end{proof}

\subsection{A lower bound for the auxiliary Lagrangian difference}
\label{sec:auxiliary-lagrangian}

The aim of this subsection is to bound from below the difference between the
value of an auxiliary Lagrangian associated with the factor $f$ and its value
at the reference factor $f_h$.  Here $f$ comes from the
decomposition $W=f\otimes f+E$, while $f_h$ is the factor defining
$W_h=f_h\otimes f_h$.  By \Cref{lem:localization-rank-one}, every
feasible graphon satisfying $I_{p_h}(W)\le I_{p_h}(W_h)$ admits such a
decomposition, with $f$ close to $u_*$ and $E$ small.  Since $0\le f\le2$,
the kernel $f\otimes f$ may take values larger than $1$ and need not be a
graphon.  We therefore introduce an auxiliary Lagrangian that can be evaluated
at both $f$ and $f_h$, temporarily leaving $E$ aside.  In
\Cref{sec:graphon-comparison}, the resulting lower bound will be
combined with estimates for $E$ to compare the actual graphons $W$ and $W_h$.

The quantities in the auxiliary Lagrangian depend only on how frequently $f$
takes each of its possible values.  Let $\lambda$ denote the Lebesgue
probability measure on $[0,1]$, and define the distributions of the values of
$f$ and $f_h$ by pushing $\lambda$ forward under these two functions:
\begin{equation*}
  \nu:=f_\#\lambda,
  \qquad
  \nu_h:=(f_h)_\#\lambda.
\end{equation*}
Thus, for every Borel set $B\subseteq[0,2]$,
$\nu(B)=\lambda\{x:f(x)\in B\}$, and the analogous identity with $f_h$ holds
for $\nu_h$.  Because $f(x)f(y)$ may exceed $1$, we use a particular auxiliary
continuation $\widetilde J_{p_h}$ of $J_{p_h}$ from $[0,1]$ to $[0,4]$, defined
below.  The first integral term in the auxiliary Lagrangian is then
\[
 \iint \widetilde J_{p_h}(f(x)f(y))\,dx\,dy.
\]
By the definition of $\nu$, the integral and the moment
$q(f)$ can be written as
\begin{equation*}
  \iint \widetilde J_{p_h}(f(x)f(y))\,dx\,dy
  =\iint \widetilde J_{p_h}(st)\,d\nu(s)d\nu(t),
  \qquad
  q(f)=\int s^d\,d\nu(s),
\end{equation*}
and the same identities hold with $f_h$ and $\nu_h$, with
$q(f_h)=q_h$.  Although $f\otimes f$ need not be a graphon, the integral
defining $t(H,\,\cdot\,)$ is still well defined for this bounded kernel.  Since
$H$ is $d$-regular, the rank-one identity gives
\begin{equation*}
  t(H,f\otimes f)=q(f)^v,
  \qquad
  t(H,W_h)=q_h^v=r_h^m.
\end{equation*}
The auxiliary Lagrangian difference studied in this subsection is
\begin{align*}
  &\left[
    \iint \widetilde J_{p_h}(f(x)f(y))\,dx\,dy
    -\mu_h\{t(H,f\otimes f)-r_h^m\}
  \right]
  \\
  &\qquad-
  \left[
    \iint \widetilde J_{p_h}(f_h(x)f_h(y))\,dx\,dy
    -\mu_h\{t(H,f_h\otimes f_h)-r_h^m\}
  \right]
  \\
  & {} = \left[
    \iint \widetilde J_{p_h}(st)\,d\nu(s)d\nu(t)
    -\mu_h\{q(f)^v-r_h^m\}
  \right]
  \\
  &\qquad\quad-
  \left[
    \iint \widetilde J_{p_h}(st)\,d\nu_h(s)d\nu_h(t)
    -\mu_h\{q_h^v-r_h^m\}
  \right].
\end{align*}
The main conclusion, \Cref{lem:auxiliary-lagrangian-bound}, provides
constants $a>0$ and $C_{d,m}\ge1$, independent of $\rho$ and $h$, with the
following property.  For every sufficiently small $\rho>0$, one can choose
$b_\rho>0$ so that, for all sufficiently small $h>0$, the above difference is
bounded below by
\begin{equation*}
  \frac{a}{2}
  \int_{\{|s-u_*|<\rho\}}(s-s_h)^2(s-t_h)^2\,d\nu(s)
  +\frac{b_\rho}{2}\nu\{|s-u_*|\ge\rho\}
  -C_{d,m}\{q(f)-q_h\}^2.
\end{equation*}
The integral in this lower bound measures the deviation, within a
neighborhood of $u_*$, from the two values $s_h,t_h$ taken by $f_h$.  The
next term measures the mass outside that neighborhood, and the final,
possibly negative, term depends only on the difference between the $d$-th
moments of $f$ and $f_h$.  This lower bound for the auxiliary Lagrangian
difference is the conclusion sought in this subsection.  It does not yet
compare the graphon costs of $W$ and $W_h$ or identify $f$ with $f_h$.  In
\Cref{sec:graphon-comparison}, we restore the remainder $E$ and
convert this bound on the auxiliary Lagrangian difference into a lower bound
involving the actual cost and constraint value of $W$.  The resulting lower
bound for the full-graphon Lagrangian difference is then used in
\Cref{sec:singular-proof} to establish optimality and uniqueness.

We now define $\widetilde J_{p_h}$.  Recall that
$\Lambda_h=\ell(p_h)-\ell_*$.  First set
\[
  \widetilde\Gamma_d(z)=
  \begin{cases}
    \Gamma_d(z),&0\le z\le1,\\
    \Gamma_d(1)+(z-1)^2,&1<z\le4,
  \end{cases}
\]
and for $0 \le z \le 4$, define
\begin{equation}\label{eq:entropy-continuation}
  \widetilde J_{p_*}(z)
  :=J_{p_*}(r_*)+\beta_d(z^d-r_*^d)
  +\widetilde\Gamma_d(z),
  \qquad
  \widetilde J_{p_h}(z)
  :=\widetilde J_{p_*}(z)+\Lambda_h z+J_{p_h}(0)-J_{p_*}(0).
\end{equation}
On $[0,1]$, this definition agrees with $J_{p_h}$ by
\eqref{eq:entropy-parameter-shift} and \eqref{eq:endpoint-supporting-gap}.  For
$1<z\le4$, however, we do not use the relative-entropy formula; instead,
\eqref{eq:entropy-continuation} means explicitly that
\[
  \widetilde J_{p_h}(z)
  =J_{p_*}(r_*)+\beta_d(z^d-r_*^d)+\Gamma_d(1)+(z-1)^2
   +\Lambda_h z+J_{p_h}(0)-J_{p_*}(0).
\]
Thus $\widetilde J_{p_h}$ is a real-valued continuous function on $[0,4]$,
but its values above $1$ are only auxiliary and are not relative entropies.
Write
$K_h(x,y):=\widetilde J_{p_h}(xy)$ for the associated kernel on $[0,2]^2$.
The following lemma gives H\"older continuity and uniform bounds for
$\widetilde J_{p_h}$ and $K_h$.
\begin{lemma}[Uniform control of auxiliary continuation and kernel]
\label{lem:continuation-kernel-bounds}
There are $C_d\ge1$ and $h_0>0$ such that, whenever $0\le h<h_0$,
$z,z'\in[0,4]$, and $x,x',y\in[0,2]$,
\begin{equation}\label{eq:continuation-kernel-bounds}
  \begin{aligned}
    |\widetilde J_{p_h}(z)-\widetilde J_{p_h}(z')|&\le C_d|z-z'|^{1/2},
    & \|\widetilde J_{p_h}-\widetilde J_{p_*}\|_{L^\infty([0,4])}&\le C_dh^2,
    & \|\widetilde J_{p_h}\|_{L^\infty([0,4])}&\le C_d,
    \\
    |K_h(x,y)-K_h(x',y)| &\le C_d|x-x'|^{1/2},
    & \|K_h-K_0\|_{L^\infty([0,2]^2)}&\le C_dh^2,
    & \|K_h\|_{L^\infty([0,2]^2)}&\le C_d.
  \end{aligned}
\end{equation}
\end{lemma}

\begin{proof}
On $[0,1]$, the only non-Lipschitz terms in
$\widetilde J_{p_*}=J_{p_*}$ are $z\log z$ and
$(1-z)\log(1-z)$.  If $0\le z'\le z\le1$,
\[
  |z\log z-z'\log z'|
  \le\int_{z'}^z|1+\log t|\,dt
  \le C_d\int_{z'}^z t^{-1/2}\,dt
  \le C_d|z-z'|^{1/2},
\]
with the same bound for $(1-z)\log(1-z)$.
On $[1,4]$, $\widetilde J_{p_*}$ is a polynomial and hence Lipschitz.
If $z\le1\le z'$, 
\[
  |\widetilde J_{p_*}(z)-\widetilde J_{p_*}(z')|
  \le |\widetilde J_{p_*}(z)-\widetilde J_{p_*}(1)|
    +|\widetilde J_{p_*}(1)-\widetilde J_{p_*}(z')|
  \le C_d\{|1-z|^{1/2}+|z'-1|^{1/2}\}
  \le C_d|z-z'|^{1/2},
\]
by continuity at $1$.
Thus $\widetilde J_{p_*}$ is $1/2$-H\"older continuous on $[0,4]$.
The definition \eqref{eq:entropy-continuation} gives
\[
  \widetilde J_{p_h}(z)-\widetilde J_{p_h}(z')
  =\widetilde J_{p_*}(z)-\widetilde J_{p_*}(z')
  +\Lambda_h(z-z'),
  \qquad
  \widetilde J_{p_h}(z)-\widetilde J_{p_*}(z)
  =\Lambda_h z+J_{p_h}(0)-J_{p_*}(0).
\]
By \eqref{eq:rank-one-parameter-expansions}, $\Lambda_h=O_d(h^2)$.
Since $p_h=(1+e^{\ell_*+\Lambda_h})^{-1}$ and
$J_p(0)=-\log(1-p)$, we also have
$J_{p_h}(0)-J_{p_*}(0)=O_d(h^2)$.
Together with $|z-z'|\le2|z-z'|^{1/2}$ on $[0,4]$, these estimates give the
desired bounds: the first identity proves the $1/2$-H\"older continuity of
$\widetilde J_{p_h}$, while the second proves the uniform bound on
$\widetilde J_{p_h}-\widetilde J_{p_*}$.
The uniform bound on $\widetilde J_{p_h}$ follows from the uniform bounds on
$\widetilde J_{p_*}$ and $\widetilde J_{p_h}-\widetilde J_{p_*}$.
The three bounds for $K_h$ follow from the corresponding bounds for
$\widetilde J_{p_h}$ by taking $z=xy$ and, for the H\"older estimate,
$z'=x'y$.
\end{proof}

The auxiliary Lagrangian comparison uses the following moment and deviation
quantities.  Since $f_h=f_{\alpha_h,s_h,t_h}$, the distribution of
its values has the explicit form
\begin{equation*}
  \nu_h=\alpha_h\delta_{s_h}+(1-\alpha_h)\delta_{t_h},
\end{equation*}
where $\delta_x$ denotes the Dirac mass at $x$.
For any probability measure $\xi$ on $[0,2]$, write $m_d(\xi)$ for its
$d$-th moment and $\Delta_h(\xi)$ for the defect of this moment relative to
$q_h$:
\begin{equation*}
  m_d(\xi):=\int x^d\,d\xi(x),
  \qquad
  \Delta_h(\xi):=m_d(\xi)-q_h.
\end{equation*}
For the two measures fixed above, $m_d(\nu)=q(f)$ and
$m_d(\nu_h)=q_h$.
In this notation, the rank-one identities above become
\begin{equation*}
  t(H,f\otimes f)=m_d(\nu)^v,
  \qquad
  t(H,W_h)=m_d(\nu_h)^v=q_h^v.
\end{equation*}
In addition to the moment defect, the lower bound controls two further ways in
which $\nu$ may differ from $\nu_h$.  Within a neighborhood of $u_*$, $\nu$
may assign mass to values other than $s_h$ and $t_h$; outside that
neighborhood, it may have nonzero mass.  To record these two deviations, for
$\rho>0$, define the central neighborhood and the tail region in $[0,2]$ by
\[
  \mathcal N_\rho:=(u_*-\rho,u_*+\rho),
  \qquad
  \mathcal T_\rho:=[0,2]\setminus\mathcal N_\rho.
\]
We always choose $\rho_0>0$ below small enough that
$\overline{\mathcal N_{\rho_0}}\subset(0,1)$.
For $h>0$, set $Q_h(x):=(x-s_h)(x-t_h)$, and for any probability measure
$\xi$ on $[0,2]$ define
\[
  A_{h,\rho}(\xi)
  :=\int_{\mathcal N_\rho}Q_h(x)^2\,d\xi(x).
\]
Here $A_{h,\rho}(\xi)$ measures deviation from the two factor values $s_h$
and $t_h$ within the central neighborhood, whereas $\xi(\mathcal T_\rho)$ is
the mass outside that neighborhood.  For the
decomposition furnished by \Cref{lem:localization-rank-one},
Chebyshev's inequality gives, for every fixed $\rho>0$, a constant
$C_{d,\rho}\ge1$ such that
\begin{equation}\label{eq:factor-tail-mass}
  \nu(\mathcal T_\rho)
  \le \rho^{-4}\|f-u_*\|_4^4
  \le C_{d,\rho}h^4.
\end{equation}

We now turn to the proof of the lower bound for the auxiliary Lagrangian
difference.  For any probability measure $\xi$ on $[0,2]$, define the
auxiliary factor functional
$\mathcal J_h$ and, using $q_h^v=r_h^m$, the auxiliary Lagrangian
$\mathcal L_h$ by
\[
  \mathcal J_h(\xi)
  :=\iint K_h(x,y)\,d\xi(x)d\xi(y),
  \qquad
  \mathcal L_h(\xi)
  :=\mathcal J_h(\xi)-\mu_h\{m_d(\xi)^v-q_h^v\}.
\]
Put $\sigma:=\xi-\nu_h$.  Expanding $\mathcal L_h(\xi)$ around $\nu_h$
produces a term linear in $\sigma$, a term quadratic in $\sigma$, and a Taylor
remainder depending on $\Delta_h(\xi)$.  The definition of $\mu_h$ and the
identity $m=dv/2$ give
$\mu_hvq_h^{v-1}=2\gamma_hq_h/d$.  To express the linear term, define the
first-variation potential $\Psi_h$ by
\begin{equation}\label{eq:first-variation}
  \Psi_h(x):=2\int K_h(x,y)\,d\nu_h(y)
  -\frac{2\gamma_hq_h}{d}x^d-\kappa_h,
\end{equation}
where $\kappa_h$ is chosen so that $\Psi_h(s_h)=0$.
Since $\int d\sigma=0$ and $\int x^d\,d\sigma=\Delta_h(\xi)$, we have
\begin{equation*}
  \mathcal J_h(\xi)-\mathcal J_h(\nu_h)
    -\mu_h\{m_d(\xi)^v-q_h^v\}
  =\int\Psi_h\,d\sigma+\iint K_h\,d\sigma d\sigma
    -\Bigl[\mu_h\{(q_h+\Delta_h(\xi))^v-q_h^v\}
      -\frac{2\gamma_hq_h}{d}\Delta_h(\xi)\Bigr].
\end{equation*}
Since $m_d(\nu_h)=q_h$, the left-hand side is precisely
$\mathcal L_h(\xi)-\mathcal L_h(\nu_h)$.  Using
$\mu_hvq_h^{v-1}=2\gamma_hq_h/d$, the bracket can be rewritten as
\[
  \mu_h\Bigl\{(q_h+\Delta_h(\xi))^v-q_h^v
  -vq_h^{v-1}\Delta_h(\xi)\Bigr\}.
\]
This is the Taylor remainder after the constant and linear terms have been
removed, so its absolute value is at most
$C_{d,m}\Delta_h(\xi)^2$.  The other two terms in the expansion are
$\int\Psi_h\,d\sigma$, which is linear in $\sigma$ and is therefore called the
first-variation term, and $\iint K_h\,d\sigma d\sigma$, which is quadratic in
$\sigma$.  We next obtain a positive lower bound for the first-variation term
and control how negative the quadratic kernel term can be.

The reference measure $\nu_h$ is supported on $\{s_h,t_h\}$.  Our choice of
$\kappa_h$ gives $\Psi_h(s_h)=0$, while the block-proportion condition
\eqref{eq:block-stationarity}, satisfied by $W_h$ by
\Cref{lem:rank-one-kkt-family}, gives $\Psi_h(t_h)=0$.  For $x$ in a
neighborhood of $u_*$, all products in the following formula lie in $(0,1)$,
so differentiating \eqref{eq:first-variation} yields
\[
  \Psi_h'(x)
  =2\left[
    \alpha_hs_hJ_{p_h}'(xs_h)
    +(1-\alpha_h)t_hJ_{p_h}'(xt_h)
  \right]-2\gamma_hq_hx^{d-1}.
\]
At $x=s_h$ and $x=t_h$, the rank-one KKT equation and
$m_d(\nu_h)=q_h$ give
\[
  \Psi_h'(x)
  =2\gamma_hq_hx^{d-1}-2\gamma_hq_hx^{d-1}
  =0 \qquad (x\in\{s_h,t_h\}).
\]
Hence $\Psi_h$ has double zeros at $s_h$ and $t_h$.
The following lemma turns these zeros into the two bounds needed below: near
$u_*$, the first variation controls the distance from $s_h$ and $t_h$, while
outside that neighborhood it imposes a fixed positive penalty.

\begin{lemma}[First variation lower bound]
\label{lem:first-variation-bound}
There exist $\rho_0>0$ and $a>0$ with the following property.  For every
$0<\rho<\rho_0$, one can choose $b_\rho,h_\rho>0$ so that, for all
$0<h<h_\rho$,
\begin{equation*}
  \Psi_h(x)\ge a Q_h(x)^2
  \quad(x\in\mathcal N_\rho),
  \qquad
  \Psi_h(x)\ge b_\rho
  \quad(x\in\mathcal T_\rho).
\end{equation*}
\end{lemma}

\begin{proof}
The central bound comes from the four equalities
$\Psi_h(s_h)=\Psi_h'(s_h)=\Psi_h(t_h)=\Psi_h'(t_h)=0$: they force
$\Psi_h$ to contain the factor $Q_h^2$, with a positive quotient near $u_*$.
The tail bound then follows from the strict positivity of the limiting
potential away from $u_*$ and uniform convergence as $h\to0$.

At $h=0$ one has $\nu_0=\delta_{u_*}$, $\gamma_0=d\beta_d$, and
$q_0=u_*^d$.  Since $\kappa_0$ is chosen so that
$\Psi_0(u_*)=0$, substituting these values into
\eqref{eq:first-variation} and using
\eqref{eq:entropy-continuation} gives
\begin{equation}\label{eq:endpoint-first-variation}
  \Psi_0(x)=2\widetilde\Gamma_d(u_*x).
\end{equation}
Near $x=u_*$, $\widetilde\Gamma_d=\Gamma_d$, and using
\eqref{eq:endpoint-gap-expansion} and
$u_*=((d-1)/d)^{1/2}$ gives
\begin{equation*}
  \Psi_0(x)
  =2\Gamma_d(r_*+u_*(x-u_*))
  =\frac{d^3}{12}(x-u_*)^4+O_d((x-u_*)^5).
\end{equation*}
To obtain a lower bound that remains uniform as the two zeros approach one
another, apply analytic division, on a neighborhood of $u_*$, to $\Psi_h(x)$
by the monic polynomial $Q_h(x)^2$.  This gives
\[
  \Psi_h(x)=Q_h(x)^2D_h(x)+S_h(x),
  \qquad \deg_xS_h\le3,
\]
where $D_h(x)$ is jointly analytic in $(h,x)$ and the coefficients of
$S_h(x)$ are analytic in $h$.  For every $h\ne0$, the four equalities above
make this remainder vanish at $s_h,t_h$, with multiplicity two;
hence $S_h(x)$ is identically zero.  By analyticity, it also vanishes at
$h=0$.  Thus
\[
  \Psi_h(x)=Q_h(x)^2D_h(x).
\]
At $h=0$ and $x=u_*$, the Taylor expansion of $\Psi_0$ gives
\[
  D_0(u_*)=\frac{d^3}{12}.
\]
By continuity, after decreasing $\rho_0$ and $h_\rho$, we therefore have
$D_h(x)\ge d^3/24$ for $0<h<h_\rho$ and $x\in\mathcal N_\rho$.  Thus the
lower bound holds with $a:=d^3/24$.

It remains to bound the tail.
By \eqref{eq:endpoint-first-variation}, $\widetilde\Gamma_d\ge0$, with equality
only at $r_*$.  Therefore, $\Psi_0\ge0$, with equality only at $u_*$.
Since $\mathcal T_\rho$ is compact and $u_* \notin \mathcal T_\rho$,
the minimum of $\Psi_0$ on $\mathcal T_\rho$ is positive.  Let $b_\rho$ be
half of this minimum:
\[
  b_\rho:=\frac12\min_{x\in\mathcal T_\rho}\Psi_0(x)>0.
\]
Writing out $\nu_h=\alpha_h\delta_{s_h}+(1-\alpha_h)\delta_{t_h}$ in
\eqref{eq:first-variation} gives
\[
  \Psi_h(x)-\Psi_0(x)
  =2\alpha_h\{K_h(x,s_h)-K_0(x,u_*)\}
  +2(1-\alpha_h)\{K_h(x,t_h)-K_0(x,u_*)\}
  -\frac{2}{d}(\gamma_hq_h-\gamma_0q_0)x^d-(\kappa_h-\kappa_0).
\]
By \Cref{lem:continuation-kernel-bounds} and the convergence
$s_h,t_h\to u_*$, we have
\[
  \|K_h(\,\cdot\,,s_h)-K_0(\,\cdot\,,u_*)\|_\infty
  +\|K_h(\,\cdot\,,t_h)-K_0(\,\cdot\,,u_*)\|_\infty
  \longrightarrow0.
\]
Also, $\alpha_h\to1/2$, $\gamma_hq_h\to\gamma_0q_0$, and
$\kappa_h\to\kappa_0$ imply $\|\Psi_h-\Psi_0\|_\infty\to0$.
Decreasing $h_\rho$ so that $\|\Psi_h-\Psi_0\|_\infty\le b_\rho$,
we get $\Psi_h\ge b_\rho$ on $\mathcal T_\rho$.
\end{proof}

Because $\Psi_h$ vanishes at both points in the support of $\nu_h$, we have
$\int\Psi_h\,d\sigma=\int\Psi_h\,d\xi$.  The preceding lemma therefore makes
the first-variation term nonnegative and quantifies its size.  The quadratic
term $\iint K_h\,d\sigma d\sigma$, however, need not be nonnegative.  The
following lemma controls its possible negative contribution when both
variables lie in $\mathcal N_\rho$.

\begin{lemma}[Central kernel bound]
\label{lem:central-kernel-bound}
Let $a>0$ be the constant from
\Cref{lem:first-variation-bound}.
There exist $\rho_0>0$ and $C_d\ge1$ with the following property.  For every
$0<\rho<\rho_0$, one can choose $h_\rho>0$ so that the following holds.  Let
$0<h<h_\rho$, let $\xi$ be any probability measure on $[0,2]$, and put
$\sigma:=\xi-\nu_h$.  Then
\begin{equation}\label{eq:central-kernel-bound}
  \iint_{\mathcal N_\rho^2}K_h\,d\sigma d\sigma
  \ge{}-\frac a2A_{h,\rho}(\xi)
  -C_d\left\{\Delta_h(\xi)^2+\xi(\mathcal T_\rho)^2\right\}.
\end{equation}
\end{lemma}

\begin{proof}
We interpolate $K_h$ in the variables $x^d$ and $y^d$ at the two reference
values $s_h^d,t_h^d$.  This separates the kernel into a part controlled by
the total mass and the $d$-th moment of $\sigma$, two parts containing one
factor of $Q_h$, and a remainder containing $Q_h(x)Q_h(y)$.  The first part
is controlled by $\Delta_h(\xi)$ and $\xi(\mathcal T_\rho)$.  The factors of
$Q_h$ allow the other three parts to be controlled by $A_{h,\rho}(\xi)$
together with these same two quantities.

To carry out this decomposition, for a function $F$ of $(x,y)$ define the
interpolation and remainder operators in the first variable by
\begin{equation}\label{eq:moment-interpolation}
  (I_h^xF)(x,y)
  :=\frac{t_h^d-x^d}{t_h^d-s_h^d}F(s_h,y)
  +\frac{x^d-s_h^d}{t_h^d-s_h^d}F(t_h,y),
  \quad
  (R_h^xF)(x,y):=F(x,y)-(I_h^xF)(x,y),
\end{equation}
and define $I_h^y,R_h^y$ analogously.  The two interpolation operators
commute.  We handle their limiting behavior at $h=0$ below using divided
differences.
Apply the exact identity
\begin{equation*}
  K_h=I_h^xI_h^yK_h+R_h^xI_h^yK_h
   +I_h^xR_h^yK_h+R_h^xR_h^yK_h.
\end{equation*}
Denote the first term by $P_h(x,y)$.  Since it is affine in each of $x^d$
and $y^d$, it can be written as
\[
  P_h(x,y)=p_{00,h}+p_{d0,h}x^d+p_{0d,h}y^d+p_{dd,h}x^dy^d.
\]
The second term vanishes at $x=s_h,t_h$ and, after
division by $Q_h(x)$, is affine in $y^d$.
The third term is the symmetric counterpart.
The last term vanishes at both $x=s_h,t_h$ and $y=s_h,t_h$.
Thus, for suitable functions $u_{0,h},u_{d,h},v_{0,h},v_{d,h}$ and
$\Theta_h$, the identity becomes
\begin{equation*}
  K_h(x,y)
  =P_h(x,y)+Q_h(x)\{u_{0,h}(x)+u_{d,h}(x)y^d\}
  +Q_h(y)\{v_{0,h}(y)+v_{d,h}(y)x^d\}
  +Q_h(x)Q_h(y)\Theta_h(x,y).
\end{equation*}

We next show that all the terms in this decomposition remain uniformly
controlled as $s_h$ and $t_h$ approach one another.  Pass to the coordinates
$w=x^d$ and $z=y^d$.
Choose $\rho_0$ and $h_\rho$ so small that
$\mathcal N_{\rho_0}$ lies
in a fixed compact interval $\mathcal U\subset(0,1)$
and $\{s_h,t_h\} \subset \mathcal N_{\rho}$.
For a function $F$ analytic on a neighborhood of $\mathcal U^2$, define
\[
  \widehat F(w,z):=F(w^{1/d},z^{1/d}),
  \qquad
  \widehat{\mathcal U}:=\{w:w^{1/d}\in\mathcal U\}.
\]
Then, $\widehat F$ is also analytic on a neighborhood of
$\widehat{\mathcal U}^2$.  In particular, its derivatives are uniformly
bounded on $\widehat{\mathcal U}^2$.
In these coordinates, $I_h^x$ and $I_h^y$ are ordinary linear interpolants
at $s_h^d,t_h^d$.  At $h=0$, these interpolation points both equal $u_*^d$, and
the interpolants converge to the corresponding tangent lines.  Thus,
although the individual coefficients in \eqref{eq:moment-interpolation}
contain $(t_h^d-s_h^d)^{-1}$, the combined interpolants and remainders have
convergent limits.  The divided-difference formulas below make this precise.

For fixed $y\in\mathcal U$, formula \eqref{eq:moment-interpolation} gives
\begin{equation*}
  \frac{R_h^xF(x,y)}{Q_h(x)}
  =\frac{x^d-s_h^d}{x-s_h}\frac{x^d-t_h^d}{x-t_h}
    \widehat F[s_h^d,x^d,t_h^d;y^d],
\end{equation*}
where the bracket denotes the second divided difference in the first
variable, with the second variable fixed at $y^d$.
The first two factors are polynomial sums and extend continuously to $x=s_h$
and $x=t_h$; for example,
\[
  \frac{x^d-s_h^d}{x-s_h}
  =\sum_{j=0}^{d-1}x^{d-1-j}s_h^j
  \le d.
\]
The divided difference has the integral representation and uniform bound:
\[
\begin{aligned}
  \widehat F[s_h^d,x^d,t_h^d;y^d]
  &=\int_0^1\int_0^{1-\theta}
    \partial_w^2\widehat F
    \!\left((1-\theta-\tau)s_h^d+\theta x^d+\tau t_h^d,y^d\right)\,d\tau\,d\theta,\\
  \left|\widehat F[s_h^d,x^d,t_h^d;y^d]\right|
  &\le \frac12\sup_{(w,z)\in\widehat{\mathcal U}^2}
    |\partial_w^2\widehat F(w,z)|.
\end{aligned}
\]
Consequently, the quotient $R_h^xF/Q_h$ extends continuously whenever any of
$\{s_h,x,t_h\}$ coincide.

We now apply the formula in both variables, taking $F=K_h$.  On
$\mathcal U^2$, all products $xy$ lie in $(0,1)$, so $K_h(x,y)=J_{p_h}(xy)$
and $K_h$ depends analytically on $(h,x,y)$.  Using brackets for the
successive second divided differences in the two variables, the last term in
the decomposition above is given exactly by
\[
  \Theta_h(x,y)
  =\prod_{c\in\{s_h,t_h\}}
    \frac{x^d-c^d}{x-c}\frac{y^d-c^d}{y-c}
    \widehat K_h[s_h^d,x^d,t_h^d;\,s_h^d,y^d,t_h^d].
\]
The successive integral representations show that this expression extends
jointly continuously when interpolation points coincide, including at
$h=0$.  The analogous lower-order divided-difference formulas and the uniform
derivative bounds on $\widehat K_h$ give a single $C_d\ge1$ such that
\[
  \max\bigl\{|p_{ij,h}|,|u_{i,h}(x)|,|v_{i,h}(y)|\bigr\}\le C_d
  \qquad  (i,j\in\{0,d\},\,0<h<h_0,\,x,y\in\mathcal N_{\rho_0}).
\]
The bounds for the one-variable functions remain valid when $x$ or $y$ equals
$s_h$ or $t_h$, by continuous extension.

The sign of the last term in the kernel decomposition is controlled by the
limiting value of $\Theta_h$.  Evaluating the jointly continuous extension at
$h=0$ and $x=y=u_*$ gives
\[
  \Theta_0(u_*,u_*)
  =\frac{(d u_*^{d-1})^4}{4}
    \left.\partial_w^2\partial_z^2 \widehat K_0(w,z)\right|_{(w,z)=(u_*^d,u_*^d)}.
\]
By \eqref{eq:entropy-continuation},
\[
  \widehat K_0(w,z)
  =\widetilde J_{p_*}((wz)^{1/d})
  =J_{p_*}(r_*)+\beta_d(wz-r_*^d)
    +\widetilde\Gamma_d\bigl((wz)^{1/d}\bigr).
\]
Since $\widetilde\Gamma_d=\Gamma_d$ near $r_*$ and
$\Gamma_d^{(j)}(r_*)=0$ for $1\le j\le3$, the chain rule leaves only the
$\Gamma_d^{(4)}(r_*)$ term, giving
\[
  \Theta_0(u_*,u_*)
  =\frac{(d u_*^{d-1})^4}{4} \left(\frac{u_*^{2-d}}{d}\right)^4 \Gamma_d^{(4)}(r_*)
  =\frac{d^3}{4}.
\]
By the joint continuity of $\Theta_h(x,y)$ at $(h,x,y)=(0,u_*,u_*)$,
after decreasing $\rho_0$ and $h_\rho$, we also have
\[
  \sup_{0<h<h_\rho}
  \left\|\Theta_h-\frac{d^3}{4}\right\|_{L^\infty(\mathcal N_\rho^2)}
  \le\frac a4.
\]

We now integrate the kernel decomposition against $\sigma\otimes\sigma$.
Since the support of $\nu_h$ lies in $\mathcal N_\rho$, we have
$\sigma(\mathcal N_\rho)=-\xi(\mathcal T_\rho)$,
$\int |Q_h|\,d\nu_h=0$, and
\begin{gather*}
  |\sigma(\mathcal N_\rho)|^2
  +\left|\int_{\mathcal N_\rho}x^d\,d\sigma(x)\right|^2
  =\xi(\mathcal T_\rho)^2
   +\left|\Delta_h(\xi)-\int_{\mathcal T_\rho}x^d\,d\xi(x)\right|^2
  \le C_d\bigl\{\Delta_h(\xi)^2+\xi(\mathcal T_\rho)^2\bigr\}, \\
  \left(\int_{\mathcal N_\rho}|Q_h|\,d|\sigma|\right)^2
  =\left(\int_{\mathcal N_\rho}|Q_h|\,d\xi\right)^2
  \le\int_{\mathcal N_\rho}Q_h^2\,d\xi
  =A_{h,\rho}(\xi).
\end{gather*}
Here the first estimate uses $x^d\le2^d$ on $[0,2]$, while the second uses
Cauchy--Schwarz.
Using these estimates together with the uniform bounds on the coefficients
and one-variable functions gives
\begin{gather*}
  \left|\iint_{\mathcal N_\rho^2}P_h\,d\sigma d\sigma\right|
  \le C_d\bigl(\Delta_h(\xi)^2+\xi(\mathcal T_\rho)^2\bigr), \\
  \left|\iint_{\mathcal N_\rho^2}Q_h(x)\{u_{0,h}(x)+u_{d,h}(x)y^d\}\,d\sigma(x)d\sigma(y)\right|
  \le \frac a8A_{h,\rho}(\xi)+C_d\bigl\{\Delta_h(\xi)^2+\xi(\mathcal T_\rho)^2\bigr\}, \\
  \left|\iint_{\mathcal N_\rho^2}Q_h(y)\{v_{0,h}(y)+v_{d,h}(y)x^d\}\,d\sigma(x)d\sigma(y)\right|
  \le \frac a8A_{h,\rho}(\xi)+C_d\bigl\{\Delta_h(\xi)^2+\xi(\mathcal T_\rho)^2\bigr\}, \\
  \iint_{\mathcal N_\rho^2}Q_h(x)Q_h(y)\Theta_h(x,y) \,d\sigma(x)d\sigma(y)
  \ge \iint_{\mathcal N_\rho^2}Q_h(x)Q_h(y) \left(\Theta_h(x,y)-\frac{d^3}{4}\right) \,d\sigma(x)d\sigma(y)
  \ge -\frac{a}{4} A_{h,\rho}(\xi).
\end{gather*}
In the last line, the omitted constant part is the nonnegative square
\[
  \frac{d^3}{4}
  \left(\int_{\mathcal N_\rho}Q_h(x)\,d\sigma(x)\right)^2.
\]
The second and third bounds use
$uv \le (a/8)u^2+(2/a)v^2$.
Combining these bounds gives \eqref{eq:central-kernel-bound}.
\end{proof}

We now combine the preceding estimates to obtain a lower bound for the
auxiliary Lagrangian difference $\mathcal L_h(\nu)-\mathcal L_h(\nu_h)$.
This is the conclusion announced at the beginning of this subsection and is
recorded in the following lemma.

\begin{lemma}[Lower bound for the auxiliary Lagrangian difference]
\label{lem:auxiliary-lagrangian-bound}
There exist $\rho_0>0$, $a>0$, and $C_{d,m}\ge1$ such that
for every $0<\rho<\rho_0$, one can choose $b_\rho,h_\rho>0$ with the following property.
Let $0<h<h_\rho$, and let $W$ satisfy the hypotheses of
\Cref{lem:localization-rank-one}.  Write $W=f\otimes f+E$ for the
decomposition furnished there, and let $\nu$ be the distribution of the
values of $f$.  Then
\begin{equation}\label{eq:auxiliary-lagrangian-bound}
\begin{aligned}
  \mathcal L_h(\nu)-\mathcal L_h(\nu_h)
  &=\mathcal J_h(\nu)-\mathcal J_h(\nu_h)
    -\mu_h\{m_d(\nu)^v-q_h^v\}\\
  &\ge \frac{a}{2} A_{h,\rho}(\nu)
    +\frac{b_\rho}{2} \nu(\mathcal T_\rho)
    -C_{d,m}\Delta_h(\nu)^2.
\end{aligned}
\end{equation}
\end{lemma}

\begin{proof}
We combine the positive first-variation bound with the lower bound for the
quadratic kernel term.  The central part follows from
\Cref{lem:central-kernel-bound} with $\xi=\nu$; the parts meeting the
tail are controlled by the uniform bounds on $K_h$ and then absorbed using the
small tail mass from \eqref{eq:factor-tail-mass}.

Put $\sigma:=\nu-\nu_h$.  Applying the expansion following
\eqref{eq:first-variation} with $\xi=\nu$, and using
$\int \Psi_h\,d\nu_h=0$ and Taylor's theorem, gives
\begin{equation*}
  \mathcal L_h(\nu)-\mathcal L_h(\nu_h)
  =\mathcal J_h(\nu)-\mathcal J_h(\nu_h)
    -\mu_h\{m_d(\nu)^v-q_h^v\}
  \ge\int\Psi_h\,d\nu+\iint K_h\,d\sigma d\sigma
    -C_{d,m}\Delta_h(\nu)^2.
\end{equation*}
The central and tail bounds in \Cref{lem:first-variation-bound}
give $\rho_0$, $a$, $b_\rho$, and $h_\rho$ such that
\[
  \int\Psi_h\,d\nu
  \ge a A_{h,\rho}(\nu)+b_\rho\nu(\mathcal T_\rho)
  \qquad
  (0<h<h_\rho,\, 0<\rho<\rho_0).
\]
Now we bound the double integral.  Split the quadratic term according to
$\mathcal N_\rho$ and $\mathcal T_\rho$:
\[
  \iint K_h\,d\sigma d\sigma
  =\iint_{\mathcal N_\rho^2}K_h\,d\sigma d\sigma
   +2\iint_{\mathcal N_\rho\times\mathcal T_\rho}K_h\,d\sigma d\sigma
   +\iint_{\mathcal T_\rho^2}K_h\,d\sigma d\sigma .
\]
The mixed rectangles have equal integrals by the symmetry of $K_h$ and
Fubini's theorem, since the continued kernel is bounded on $[0,2]^2$.
The central-kernel bound \eqref{eq:central-kernel-bound} controls the
first integral.  It remains to control the two parts meeting
$\mathcal T_\rho$.
After decreasing $h_\rho$, we have
$\{s_h,t_h\}\subset\mathcal N_\rho$, and hence
$\sigma(\mathcal N_\rho)=-\nu(\mathcal T_\rho)$ and $\sigma=\nu$ on
$\mathcal T_\rho$.  Thus, for $y\in\mathcal T_\rho$,
\[
\begin{aligned}
  \int_{\mathcal N_\rho}K_h(x,y)\,d\sigma(x)
  &=\int_{\mathcal N_\rho}
    \{K_h(x,y)-K_h(u_h,y)\}\,d\sigma(x)
    -\nu(\mathcal T_\rho) K_h(u_h,y).
\end{aligned}
\]
Since $|\sigma|\le\nu+\nu_h$, the H\"older bound in
\Cref{lem:continuation-kernel-bounds}, together with
\eqref{eq:rank-one-reduction-bounds} and
\eqref{eq:rank-one-parameter-expansions}, gives
\begin{equation*}
\begin{aligned}
  \left|\int_{\mathcal N_\rho}
    \{K_h(x,y)-K_h(u_h,y)\}\,d\sigma(x)\right|
  \le C_d\int_{\mathcal N_\rho}|x-u_h|^{1/2}\,d(\nu+\nu_h)(x)
  \le C_d\left\{\|f-u_h\|_4^{1/2}+h^{1/2}\right\}
  \le C_dh^{1/2}.
\end{aligned}
\end{equation*}
The uniform bound on $K_h$ from
\Cref{lem:continuation-kernel-bounds} therefore gives
\[
  \left|\int_{\mathcal N_\rho}K_h(x,y)\,d\sigma(x)\right|
  \le C_dh^{1/2}+C_d\nu(\mathcal T_\rho).
\]
Integrating over $y\in\mathcal T_\rho$ and using the same uniform bound on
$K_h$ gives the two remaining estimates:
\[
  \left|2\iint_{\mathcal N_\rho\times\mathcal T_\rho} K_h\,d\sigma d\sigma\right|
  \le C_dh^{1/2}\nu(\mathcal T_\rho)+C_d\nu(\mathcal T_\rho)^2,
  \qquad
  \left|\iint_{\mathcal T_\rho^2}K_h\,d\sigma d\sigma\right|
  \le C_d\nu(\mathcal T_\rho)^2.
\]
Combining the first-variation, central, mixed, and tail bounds gives
\begin{equation*}
  \mathcal J_h(\nu)-\mathcal J_h(\nu_h)
    -\mu_h\{m_d(\nu)^v-q_h^v\}
  \ge\frac{a}{2}A_{h,\rho}(\nu)+b_\rho\nu(\mathcal T_\rho)
  -C_{d,m}\Delta_h(\nu)^2-C_d\nu(\mathcal T_\rho)^2
    -C_dh^{1/2}\nu(\mathcal T_\rho).
\end{equation*}
Using the tail bound \eqref{eq:factor-tail-mass}, decrease $h_\rho$ so that
$C_dh^{1/2}\le b_\rho/4$ and $C_d\nu(\mathcal T_\rho)\le b_\rho/4$.
Then, \eqref{eq:auxiliary-lagrangian-bound} follows.
\end{proof}

\subsection{The arbitrary-graphon comparison}
\label{sec:graphon-comparison}

We now restore the remainder $E$ and convert the lower bound for the auxiliary
Lagrangian difference into the corresponding bound for the Lagrangian
difference between the actual graphons $W$ and $W_h$.
By \Cref{lem:localization-rank-one}, every feasible graphon satisfying
$I_{p_h}(W)\le I_{p_h}(W_h)$ can be written as $W=f\otimes f+E$, while
\Cref{lem:auxiliary-lagrangian-bound} controls the auxiliary Lagrangian
difference between $f$ and $f_h$.  The next lemma estimates the contribution
of $E$ and gives the required full-graphon bound.

\begin{lemma}[Lower bound for the full-graphon Lagrangian difference]
\label{lem:graphon-lagrangian-bound}
There exist $\rho_0>0$ and $C_{d,m}\ge1$ such that for
every $0<\rho<\rho_0$, one can choose $C_{d,\rho}\ge1$ and $h_\rho>0$ with
the following property.  Let $0<h<h_\rho$, and let $W$ satisfy the hypotheses of
\Cref{lem:localization-rank-one}.  Write $W=f\otimes f+E$ for the
decomposition provided there, and let $\nu$ be the distribution of the
values of $f$.  Then,
\begin{equation}\label{eq:graphon-lagrangian-bound}
\begin{aligned}
  &\left[I_{p_h}(W)-\mu_h\{t(H,W)-r_h^m\}\right]
   -\left[I_{p_h}(W_h)-\mu_h\{t(H,W_h)-r_h^m\}\right]\\
  &\qquad\qquad\qquad\qquad\qquad\qquad\qquad=I_{p_h}(W)-I_{p_h}(W_h)-\mu_h\{t(H,W)-r_h^m\}\\
  &\qquad\qquad\qquad\qquad\qquad\qquad\qquad
  \ge C_{d,\rho}^{-1}\!\left[A_{h,\rho}(\nu)+\nu(\mathcal T_\rho)+\|E\|_2^2\right]
  -C_{d,m}\Delta_h(\nu)^2.
\end{aligned}
\end{equation}
\end{lemma}

\begin{proof}
Put $q:=q(f)=m_d(\nu)$ and
$\Omega_\rho:=\{x:f(x)\in\mathcal N_\rho\}$.  Then
$\lambda(\Omega_\rho^c)=\nu(\mathcal T_\rho)$.
Combining the localization results in \Cref{lem:localization-rank-one}
and \eqref{eq:factor-tail-mass}, we may use throughout the proof
\begin{equation}\label{eq:graphon-comparison-estimates}
  \begin{gathered}
    \nu(\mathcal T_\rho)\le C_{d,\rho}h^4,
    \qquad |\Delta_h(\nu)|\le C_dh^2,
    \qquad \|E\|_2\le C_dh,
    \qquad \|f-u_*\|_4\le C_dh,\\
    |t(H,W)-q^v|\le C_{d,m}\|E\|_2^3,
    \qquad \lambda(\Omega_\rho)\ge\frac12,
    \qquad 0\le f\le2,
    \qquad \|E\|_\infty\le5.
  \end{gathered}
\end{equation}
The bound on $\lambda(\Omega_\rho)$ follows from
$\lambda(\Omega_\rho^c)=\nu(\mathcal T_\rho)$, while the last bound uses
$0\le W\le1$ and $0\le f\le2$.
Since
\[
  \mathcal J_h(\nu)
  =\iint\widetilde J_{p_h}(f(x)f(y))\,dx\,dy,
  \qquad
  \mathcal J_h(\nu_h)=I_{p_h}(W_h),
\]
and $t(H,W_h)=q_h^v=r_h^m$, the full-graphon Lagrangian difference has the
exact decomposition
\begin{equation*}
\begin{aligned}
  &I_{p_h}(W)-I_{p_h}(W_h)-\mu_h\{t(H,W)-r_h^m\}\\
  &\quad=\bigl[\mathcal L_h(\nu)-\mathcal L_h(\nu_h)\bigr]
  +\iint\{J_{p_h}(W(x,y))-\widetilde J_{p_h}(f(x)f(y))\}\,dx\,dy\\
  &\qquad-\mu_h\{t(H,W)-q^v\}.
\end{aligned}
\end{equation*}
The first term is controlled by \Cref{lem:auxiliary-lagrangian-bound}, and
the last by the homomorphism-density estimate in
\eqref{eq:rank-one-reduction-bounds}.  It remains to bound the middle
double integral over its central, mixed, and tail--tail parts.

We first consider the central square $\Omega_\rho^2$.  Decrease $\rho_0$ so
that $f(x)f(y)$ lies in a fixed closed subinterval of $(0,1)$ whenever
$x,y\in\Omega_\rho$.  Put
\[
  \|E\|_{2,G}
  :=\left( \iint_G E(x,y)^2\,dx\,dy \right)^{1/2}
  \qquad (G \subset[0,1]^2).
\]
On $\Omega_\rho^2$, the identity $\mathcal M_h'(xy)=J_{p_h}'(xy)-\gamma_h(xy)^{d-1}$ yields
the decomposition of the integrand:
\begin{equation*}
  J_{p_h}(W)-J_{p_h}(f\otimes f)
  = \bigl[ J_{p_h}(W)-J_{p_h}(f\otimes f) -(J_{p_h})'(f\otimes f)E \bigr]
  +\mathcal M_h'(f\otimes f)E
      +\gamma_h(f\otimes f)^{d-1}E.
\end{equation*}
In this step, we bound the integrals of these three terms on $\Omega_\rho^2$ separately.
Since $J_{p_h}''(z)=1/(z(1-z))\ge4$, strong convexity gives the first bound
\[
  \iint_{\Omega_\rho^2}
  \{J_{p_h}(W)-J_{p_h}(f\otimes f)-(J_{p_h})'(f\otimes f)E\}
  \ge2\|E\|_{2,\Omega_\rho^2}^2.
\]
For the second bound, we divide $\mathcal M_h'$ by $Q_h$ in each variable and write
\[
  \mathcal M_h'(xy)=Q_h(x)U_h(x,y)+Q_h(y)V_h(x,y)+C_h(x,y),
\]
where $C_h$ has degree at most one in each variable.  The factorization
\eqref{eq:mh-derivative-factorization} then gives
$\mathcal M_h'(xy)=0$ for every $(x,y)\in\{s_h,t_h\}^2$.  Since
$Q_h(s_h)=Q_h(t_h)=0$, the bilinear remainder $C_h$ must vanish identically.
At $h=0$, $\mathcal M_0'(xy)$ is an analytic multiple of $(xy-r_*)^3$ and
$Q_0(z)=(z-u_*)^2$, so $U_0$ and $V_0$ vanish at $(u_*,u_*)$.  By
joint continuity, after decreasing $\rho_0$ and $h_\rho$, we may ensure that
$|U_h|$ and $|V_h|$ are at most $\sqrt{a/4}$ on $\mathcal N_\rho^2$, and
\[
  |\mathcal M_h'(xy)|
  \le\sqrt{\frac a4}
    \{|Q_h(x)|+|Q_h(y)|\}
  \qquad(x,y\in\mathcal N_\rho,\ 0<h<h_\rho).
\]
Then, $uv \le (1/4)u^2+v^2$ and
$(u+v)^2\le 2(u^2+v^2)$ give
\begin{align*}
  \left|\iint_{\Omega_\rho^2} \{\mathcal M_h'(f(x)f(y))\} E(x,y)\,dx\,dy\right|
  &\le \frac{1}{4} \iint_{\Omega_\rho^2} \frac{a}{4}\{|Q_h(f(x))|+|Q_h(f(y))|\}^2 \,dx\,dy + \|E\|_{2,\Omega_\rho^2}^2 \\
  &\le \frac{a}{16} \iint_{\Omega_\rho^2} 2|Q_h(f(x))|^2+2|Q_h(f(y))|^2 \,dx\,dy + \|E\|_{2,\Omega_\rho^2}^2 \\
  &\le \frac a4 A_{h,\rho}(\nu) + \|E\|_{2,\Omega_\rho^2}^2.
\end{align*}
For the third bound, the orthogonality in \eqref{eq:rank-one-orthogonality} moves
the remaining linear term to the tail--tail square:
\[
  \left| \iint_{\Omega_\rho^2}\gamma_h(f(x)f(y))^{d-1}E(x,y)\,dx\,dy \right|
  = \left| \iint_{(\Omega_\rho^c)^2}\gamma_h(f(x)f(y))^{d-1}E(x,y)\,dx\,dy \right|
  \le C_d \nu(\mathcal T_\rho)^2.
\]
Combining the three bounds gives
\begin{equation*}
  \iint_{\Omega_\rho^2} \{J_{p_h}(W)-J_{p_h}(f\otimes f)\}
  \ge \|E\|_{2,\Omega_\rho^2}^2 -\frac a4 A_{h,\rho}(\nu)-C_d\nu(\mathcal T_\rho)^2.
\end{equation*}

We next consider the mixed rectangles.  For almost every $x$, the
orthogonality in
\eqref{eq:rank-one-orthogonality}, together with
\eqref{eq:graphon-comparison-estimates}, gives
\begin{equation*}
  u_*^{d-1}\left|\int_{\Omega_\rho}E(x,y)\,dy\right|
  \le\int_{\Omega_\rho}
    |f(y)^{d-1}-u_*^{d-1}|\,|E(x,y)|\,dy
  +\int_{\Omega_\rho^c}f(y)^{d-1}|E(x,y)|\,dy
  \le C_{d,\rho}h.
\end{equation*}
Consequently, if $\overline W_x:= (\int_{\Omega_\rho}W(x,y)\,dy)/\lambda(\Omega_\rho)$,
then
\begin{equation*}
  |\overline W_x-u_*f(x)|
  \le |f(x)|\left|
    \frac1{\lambda(\Omega_\rho)}\int_{\Omega_\rho} (f(y)-u_*)\,dy\right|
    +\frac1{\lambda(\Omega_\rho)}\left|\int_{\Omega_\rho}E(x,y)\,dy\right|
  \le C_{d,\rho}h.
\end{equation*}
Since $\overline W_x\in[0,1]$, we have
$J_{p_h}(\overline W_x)=\widetilde J_{p_h}(\overline W_x)$.
By Jensen's and H\"older's inequalities, with
\eqref{eq:continuation-kernel-bounds} and
\eqref{eq:graphon-comparison-estimates},
\begin{equation*}
\begin{aligned}
&\int_{\Omega_\rho}
  \{J_{p_h}(W(x,y))-\widetilde J_{p_h}(f(x)f(y))\}\,dy\\
&\quad\ge \lambda(\Omega_\rho) \left\{J_{p_h}(\overline W_x)-\widetilde J_{p_h}(u_*f(x))\right\}
-\int_{\Omega_\rho} \left\{\widetilde J_{p_h}(f(x)f(y))-\widetilde J_{p_h}(u_*f(x))\right\}\,dy\\
&\quad\ge-C_d |\overline W_x - u_*f(x)|^{1/2} -C_d\int_{\Omega_\rho}|f(y)-u_*|^{1/2}\,dy \ge-C_{d,\rho}h^{1/2} \ge -\frac{b_\rho}{16}.
\end{aligned}
\end{equation*}
Integrating this bound over $x\in\Omega_\rho^c$ yields
\begin{equation*}
\begin{aligned}
  2\iint_{\Omega_\rho^c\times \Omega_\rho}
  \{J_{p_h}(W)-\widetilde J_{p_h}(f\otimes f)\}
  \ge -\frac{b_\rho}{8}\nu(\mathcal T_\rho).
\end{aligned}
\end{equation*}

It remains to consider the tail--tail square.  The uniform bound in
\eqref{eq:continuation-kernel-bounds} gives
\[
  \iint_{(\Omega_\rho^c)^2}
  \{J_{p_h}(W)-\widetilde J_{p_h}(f\otimes f)\}
  \ge-C_d\nu(\mathcal T_\rho)^2.
\]

Combining all bounds with \eqref{eq:auxiliary-lagrangian-bound}
and $|t(H,W)-q^v|\le C_{d,m}\|E\|_2^3$ gives
\begin{equation*}
\begin{aligned}
  I_{p_h}(W)-I_{p_h}(W_h)-\mu_h\{t(H,W)-r_h^m\}
  &\ge\frac a4 A_{h,\rho}(\nu)
    +\frac{3b_\rho}{8}\nu(\mathcal T_\rho)
    +\|E\|_{2,\Omega_\rho^2}^2\\
  &\quad-C_{d,m}\Delta_h(\nu)^2
    -C_d\nu(\mathcal T_\rho)^2
    -C_{d,m}\|E\|_2^3.
\end{aligned}
\end{equation*}
Since $\|E\|_\infty\le5$ and
$\lambda([0,1]^2\setminus\Omega_\rho^2)\le2\nu(\mathcal T_\rho)$, choose
$C_{d,\rho}\ge1$ so that $50C_{d,\rho}^{-1}\le b_\rho/16$.  Then
\begin{equation*}
  \|E\|_{2,\Omega_\rho^2}^2
  \ge C_{d,\rho}^{-1}\|E\|_{2,\Omega_\rho^2}^2
  =C_{d,\rho}^{-1}\left(
    \|E\|_2^2-\|E\|_{2,[0,1]^2\setminus\Omega_\rho^2}^2\right)
  \ge C_{d,\rho}^{-1}\|E\|_2^2
    -50C_{d,\rho}^{-1}\nu(\mathcal T_\rho)
  \ge C_{d,\rho}^{-1}\|E\|_2^2
    -\frac{b_\rho}{16}\nu(\mathcal T_\rho).
\end{equation*}
By \eqref{eq:graphon-comparison-estimates}, we may also decrease $h_\rho$ so that
$C_d\nu(\mathcal T_\rho)\le b_\rho/16$ and $C_{d,m}\|E\|_2\le C_{d,\rho}^{-1}/2$.
Consequently,
\begin{equation*}
  I_{p_h}(W)-I_{p_h}(W_h)-\mu_h\{t(H,W)-r_h^m\}
  \ge\frac a4 A_{h,\rho}(\nu)
  +\frac{b_\rho}{4}\nu(\mathcal T_\rho)
  +\frac{1}{2}C_{d,\rho}^{-1}\|E\|_2^2
  -C_{d,m}\Delta_h(\nu)^2.
\end{equation*}
Increasing $C_{d,\rho}$ if necessary proves \eqref{eq:graphon-lagrangian-bound}.
\end{proof}

\subsection{Proof of the singular-endpoint theorem}
\label{sec:singular-proof}

We finish by applying the lower bound for the full-graphon Lagrangian
difference to the decomposition supplied by the localization and rank-one
reduction.

\begin{proof}[Proof of \Cref{thm:singular-endpoint}]
\Cref{lem:rank-one-kkt-family} provides the real-analytic map
$h\mapsto(p_h,r_h,u_h,\alpha_h)$ for $|h|<h_0$, with the desired
values at $h=0$.  For every measurable $A_h\subset[0,1]$ of measure
$\alpha_h$, the corresponding factor is $f_{\alpha_h,s_h,t_h}$, up to measure-preserving relabeling.
Since $t_h-s_h=2h\ne0$ for $h>0$, the resulting graphon $W_h$ is
nonconstant, rank-one, and bipodal, and the same lemma gives
$t(H,W_h)=r_h^m$.  Its comparison with the constant graphon and
the inequality $p_h<\pc(r_h)$ follow from \Cref{lem:constant-graphon-comparison}.

It remains to prove optimality and uniqueness.
Fix $0<\rho<\rho_0$ and $0<h<h_\rho$ as in \Cref{lem:graphon-lagrangian-bound}.
Assume $I_{p_h}(W)\le I_{p_h}(W_h)$ and $W$ is feasible.
By \Cref{lem:localization-rank-one}, write
\[
  W=f\otimes f+E,
  \qquad
  T_E(f^{d-1})=0,
  \qquad
  q:=q(f).
\]
Let $\nu$ be the distribution of the values of $f$, so that $q=m_d(\nu)$ and
$\Delta_h(\nu)=q-q_h$.  Applying \Cref{lem:graphon-lagrangian-bound} gives
\eqref{eq:graphon-lagrangian-bound}.  Its left-hand side is nonpositive:
\[
  I_{p_h}(W)-I_{p_h}(W_h)-\mu_h\{t(H,W)-r_h^m\}
  \le I_{p_h}(W)-I_{p_h}(W_h)\le0,
\]
because $W$ is feasible and $\mu_h>0$.  Comparing this upper bound with the
lower bound in \eqref{eq:graphon-lagrangian-bound}, and increasing the
constant if necessary, gives $C_{d,m,\rho}\ge1$ such that
\[
  A_{h,\rho}(\nu)+\nu(\mathcal T_\rho)+\|E\|_2^2
  \le C_{d,m,\rho}^2\Delta_h(\nu)^2,
  \qquad
  \|E\|_2\le C_{d,m,\rho}|\Delta_h(\nu)|.
\]
Combining the homomorphism-density bound in
\eqref{eq:rank-one-reduction-bounds} with Taylor's theorem gives
\begin{equation*}
\begin{aligned}
  \left|t(H,W)-r_h^m-vq_h^{v-1}\Delta_h(\nu)\right|
  \le 
  |t(H,W)-q^v|
    +\left|(q_h+\Delta_h(\nu))^v-q_h^v
      -vq_h^{v-1}\Delta_h(\nu)\right|
  \le C_{d,m,\rho}\Delta_h(\nu)^2.
\end{aligned}
\end{equation*}
Since $q_h$ stays bounded away from zero, the linear term in the preceding
estimate dominates the quadratic error when $\Delta_h(\nu)<0$.  Feasibility
would then be violated, so, after decreasing $h_\rho$, we have
$\Delta_h(\nu)\ge0$.  The localization estimate
\eqref{eq:rank-one-reduction-bounds}, together with the positive lower
bounds for $q_h$ and $\mu_h$, allows us to decrease $h_\rho$ further so that
\[
  C_{d,m,\rho}\Delta_h(\nu)\le\frac14vq_h^{v-1},
  \qquad
  C_{d,m,\rho}\Delta_h(\nu)\le\frac14\mu_hvq_h^{v-1}.
\]
This yields
\begin{equation*}
  \mu_h\{t(H,W)-r_h^m\}-C_{d,m,\rho}\Delta_h(\nu)^2
  \ge\frac12\mu_hvq_h^{v-1}\Delta_h(\nu).
\end{equation*}
Substituting this into \eqref{eq:graphon-lagrangian-bound} gives
\begin{equation*}
  I_{p_h}(W)-I_{p_h}(W_h)
  \ge
  \left(\frac12\mu_hvq_h^{v-1}\Delta_h(\nu)
    +C_{d,m,\rho}^{-1}
    \left[A_{h,\rho}(\nu)+\nu(\mathcal T_\rho)+\|E\|_2^2\right]\right)
  \ge0.
\end{equation*}
The left-hand side is nonpositive by the upper bound above, so equality holds
throughout.  Consequently,
\begin{equation*}
  \Delta_h(\nu)=0,
  \qquad A_{h,\rho}(\nu)=0,
  \qquad \nu(\mathcal T_\rho)=0,
  \qquad E=0\quad\text{a.e.}
\end{equation*}
Since $A_{h,\rho}(\nu)=0$ and $\nu(\mathcal T_\rho)=0$,
$\nu$ is supported on $\{s_h,t_h\}$, and $\Delta_h(\nu)=0$ implies
$\nu=\alpha_h\delta_{s_h}+(1-\alpha_h)\delta_{t_h}$.
Moreover, $E\equiv0$ implies that $W=f\otimes f$ and $W$ is a measure-preserving relabeling of $W_h$.
\end{proof}

